\documentclass[11pt,a4paper]{article}
\usepackage[utf8]{inputenc}
\usepackage[T1]{fontenc}
\usepackage[spanish,es-noshorthands]{babel}
\usepackage[margin=2.5cm]{geometry}
\usepackage{parskip}
\usepackage{hyperref}
\usepackage{enumitem}
\usepackage{booktabs}
\usepackage{multirow}
\usepackage{array}
\usepackage{xcolor}
\usepackage{graphicx}
\usepackage{pgfplots}
\pgfplotsset{compat=1.18}
\usetikzlibrary{arrows.meta,positioning,fit,shapes.multipart}
\usepackage{longtable}
\usepackage{amsmath}
\usepackage{amssymb}

\definecolor{done}{RGB}{0,140,0}
\definecolor{partialcol}{RGB}{200,120,0}
\definecolor{pendingcol}{RGB}{180,0,0}
\definecolor{bugcol}{RGB}{140,0,140}

\setlength{\parindent}{0pt}
\setlength{\parskip}{6pt}

\hypersetup{colorlinks=true, linkcolor=black, urlcolor=blue}

\title{\textbf{Búsqueda Semántica de Objetos en Simulación\\con Razonamiento por Habitaciones}\\[6pt]
\large Estado del proyecto y plan de desarrollo\\[4pt]
\normalsize VLMaps + YOLOE + Habitat-Sim \,$\cdot$\, HSSD como dataset principal\\
Extensión futura: ROS\,1 Noetic + Toyota HSR}
\author{Mario García Blasco}
\date{22 de abril de 2026}

\begin{document}

\maketitle
\tableofcontents
\newpage

% =============================================================================
\section{Propósito del proyecto}
% =============================================================================

\subsection{Problema}

Dado un objetivo expresado en lenguaje natural ---por ejemplo, \textit{``encuentra
la botella de agua''}, \textit{``busca la taza''} o \textit{``ve a la cocina y
mira si hay un portátil en la mesa''}---, un robot simulado debe localizar el
objeto de forma eficiente en un entorno doméstico interior.

El robot opera sobre escenas HSSD (\textit{Habitat Synthetic Scenes Dataset})
renderizadas con Habitat-Sim. La escena principal es \textbf{102344280}, una
vivienda de una planta con 15~habitaciones anotadas ($\sim$447\,m$^2$). El robot
no conoce la posición exacta de los objetos: dispone únicamente de un mapa
semántico pre-construido (VLMap), de las anotaciones nativas de habitaciones de
HSSD y de la imagen RGB de su cámara.

El objetivo del proyecto no es ampliar infraestructura genérica de simulación ni
acumular datasets, sino construir un sistema de búsqueda que entienda la
estructura semántica de la vivienda. El robot debe razonar sobre dónde tiene
sentido buscar antes de lanzarse a inspeccionar candidatos arbitrarios, y debe
ser capaz de encontrar incluso objetos que ni siquiera están explícitamente
representados en el mapa semántico (botellas, tazas, portátiles, cajas, etc.).

\subsection{Hipótesis}

La hipótesis central del TFG es la siguiente:

\begin{quote}
\textit{Un robot que razona explícitamente sobre la estructura de la casa
---primero por habitación, después por mueble o zona y finalmente por
verificación visual del objeto--- encuentra objetos domésticos de forma más
eficiente que un sistema que busca sobre un heatmap semántico global o que se
apoya en un único nivel de información.}
\end{quote}

«Más eficiente» se traduce en métricas concretas: menos habitaciones visitadas
antes del éxito, menos candidatos inspeccionados, mayor proporción de aciertos
al primer intento y caminos más cortos. Estas métricas se introducen formalmente
en la Sección~\ref{sec:fase-h}.

\subsection{Tres niveles de conocimiento}
\label{sec:tres-niveles}

El sistema combina tres niveles semánticos jerárquicos. Cada uno aporta
información que ningún otro puede sustituir:

\begin{center}
\begin{tabular}{p{2.8cm}p{4.5cm}p{6.5cm}}
\toprule
\textbf{Nivel} & \textbf{Fuente} & \textbf{Qué aporta} \\
\midrule
1.~Habitación &
  \texttt{semantic\_config.json} de HSSD &
  «Esto es la cocina; aquello es el baño». Polígonos de planta con etiqueta. \\
2.~Mueble / zona &
  VLMap (36 categorías HSSD) &
  «Aquí hay una encimera; ahí una mesa». Localización aproximada de las
  superficies sobre las que típicamente reposan los objetos. \\
3.~Objeto específico &
  YOLOE (vocabulario abierto) &
  «Sobre esa encimera hay una botella». Verificación visual final con la
  cámara del robot. \\
\bottomrule
\end{tabular}
\end{center}

\subsection{Por qué un solo nivel no basta}

Cualquiera de los tres niveles utilizado de forma aislada produce un sistema
incompleto:

\begin{itemize}
  \item \textbf{VLMaps por sí solo no encuentra objetos pequeños}. Las
        categorías indexadas en el VLMap son muebles y superficies estáticas
        (\texttt{table}, \texttt{counter}, \texttt{sofa}, \texttt{bed},
        \texttt{cabinet}, etc.). Objetos de interés cotidianos como
        \texttt{bottle}, \texttt{cup}, \texttt{box} o \texttt{laptop} no
        aparecen en el VLMap, porque CLIP los confunde fácilmente con los
        muebles que los soportan o porque su huella es demasiado pequeña en la
        proyección 2D. Sin una capa adicional, el sistema no tiene forma de
        \emph{encontrar} esos objetos: como mucho, puede localizar el mueble
        donde típicamente residen.
  \item \textbf{YOLOE por sí solo no sabe a dónde ir}. YOLOE es un detector de
        vocabulario abierto que opera sobre la imagen actual. No tiene memoria
        espacial ni mapa, por lo que sin un mecanismo que lo lleve frente a las
        zonas plausibles de la casa, su trabajo se reduce a inspeccionar el
        campo de visión actual. Una búsqueda basada exclusivamente en YOLOE
        degenera en exploración ciega.
  \item \textbf{El razonamiento por habitaciones es el pegamento que falta}.
        Saber que una botella vive típicamente en la cocina permite al robot
        saltarse cuatro habitaciones que no aportan información. Saber que
        dentro de la cocina debe mirar primero encimeras y mesas reduce
        drásticamente el número de candidatos a verificar con YOLOE. El
        razonamiento por habitaciones convierte una búsqueda exhaustiva en una
        búsqueda dirigida.
\end{itemize}

\subsection{Idea clave: razonar antes de moverse}

La política de búsqueda del sistema sigue siempre la misma estructura
jerárquica:

\begin{enumerate}
  \item Identificar las habitaciones de la escena y construir un \emph{prior}
        sobre cuál es la más probable para el objeto buscado.
  \item Traducir el objeto a un conjunto de categorías VLMap plausibles, de
        forma que el sistema sepa qué muebles o zonas inspeccionar dentro de la
        habitación elegida (por ejemplo \texttt{laptop} $\rightarrow$
        \texttt{[table, counter, shelving]}).
  \item Localizar instancias de esas categorías \emph{dentro de la habitación
        seleccionada}, no globalmente.
  \item Navegar al candidato más prometedor.
  \item Verificar visualmente con YOLOE si el objeto está allí.
  \item Si la verificación falla, continuar con el siguiente candidato dentro
        de la misma habitación; si se agotan, pasar a la siguiente habitación
        más probable. Si la verificación tiene éxito, declarar el objeto
        \textbf{ENCONTRADO}.
\end{enumerate}

Este flujo es la traducción operativa de la hipótesis. La diferencia decisiva
con respecto a un sistema de búsqueda plana está en el orden: se decide
\emph{en qué habitación} buscar antes de decidir \emph{qué candidato}
inspeccionar.

\subsection{De búsqueda plana (M2) a búsqueda jerárquica (room-aware)}
\label{sec:m2-vs-room-aware}

El sistema dispone hoy de la infraestructura necesaria para ambos estilos de
búsqueda. La diferencia conceptual entre ellos es la siguiente:

\begin{itemize}
  \item \textbf{Búsqueda plana (M2)}. Para cada categoría se construye un
        heatmap global sobre todo el mapa, se extraen componentes conexos, se
        ordenan por una métrica de calidad y se navega al primero. La
        habitación a la que pertenece el candidato es información incidental:
        no participa en la decisión. El robot puede cruzar la casa entera para
        ir a un candidato ruidoso situado en una habitación irrelevante.
  \item \textbf{Búsqueda jerárquica (room-aware)}. La decisión se descompone en
        dos niveles. Primero se elige la habitación con mayor utilidad esperada
        para la query (combinando priors LLM, evidencia del heatmap, distancia
        y exploración previa). Después se ejecuta el pipeline M2, pero
        restringiendo los candidatos a los que caen dentro de esa habitación.
        Si la habitación elegida se agota sin éxito, se pasa a la siguiente
        más probable.
\end{itemize}

El sistema room-aware no reemplaza al pipeline M2: lo \emph{encapsula}. Toda la
maquinaria de heatmaps, planificación y verificación YOLOE sigue siendo
relevante; la diferencia es que ahora actúa dentro de la habitación correcta y
en el orden correcto. La Sección~\ref{sec:room-aware} desarrolla el plan
concreto para cerrar este bucle.

% =============================================================================
\section{Alcance y decisiones de diseño}
% =============================================================================

\subsection{HSSD como dataset principal}

HSSD es un conjunto de escenas de interiores sintéticas de alta calidad
desarrolladas para Habitat-Sim. Su elección como dataset principal responde a
tres motivos directamente alineados con la hipótesis del proyecto:

\begin{itemize}
  \item \textbf{Habitaciones anotadas nativamente}. Cada escena incluye un
        fichero \texttt{semantics/scenes/<id>.semantic\_config.json} con
        polígonos de planta y nombres de habitación. La clase
        \texttt{SemanticSceneRoomProvider} los rasteriza directamente sobre el
        grid VLMap sin intervención manual. El razonamiento por habitaciones
        descansa sobre estas anotaciones, por lo que su disponibilidad nativa
        es decisiva.
  \item \textbf{Geometría limpia e iluminación controlada}. Las escenas
        sintéticas facilitan la verificación visual con YOLOE y eliminan
        artefactos de reconstrucción fotogramétrica que afectan a otros
        datasets.
  \item \textbf{Alineación con el objetivo del proyecto}. Sin anotaciones de
        habitación fiables, el proyecto se vería forzado a etiquetar
        manualmente cada escena, desplazando esfuerzo del problema central
        (razonamiento jerárquico) a un problema instrumental (anotación).
\end{itemize}

La escena principal es \textbf{102344280}: una vivienda de una planta con
15~habitaciones (\texttt{kitchen}, \texttt{dining room}, \texttt{bedroom}$\times$3,
\texttt{bathroom}$\times$3, \texttt{hallway}, \texttt{entryway},
\texttt{closet}$\times$3, \texttt{tv room}).

\subsection{MP3D como soporte secundario}

Toda la funcionalidad MP3D existente se mantiene intacta: el sistema sigue
funcionando sobre escenas MP3D y existe un \texttt{LabelMeRoomProvider} listo
para usar si se ejecuta el etiquetado manual. Sin embargo, MP3D no condiciona
la hoja de ruta actual:

\begin{itemize}
  \item El etiquetado LabelMe de las 3~escenas MP3D queda \textbf{diferido}
        como tarea secundaria.
  \item La construcción de VLMaps para escenas MP3D adicionales queda diferida.
  \item Las limitaciones visuales de MP3D (texturas ruidosas, reconstrucciones
        imperfectas) perjudican la verificación YOLOE y no son representativas
        de la capacidad real del método.
\end{itemize}

\subsection{El razonamiento por habitaciones es la prioridad}

La siguiente fase del proyecto es que el robot pueda recibir instrucciones del
estilo \textit{``encuentra la taza''}, \textit{``ve a la cocina y mira si hay
una botella en la encimera''} o \textit{``busca el portátil''}, y resuelva la
búsqueda usando los tres niveles semánticos descritos en la
Sección~\ref{sec:tres-niveles}.

Esta fase es prioritaria frente a la inserción de nuevos objetos en
simulación, los baselines M1/M3 y la evaluación cuantitativa amplia. Sin
razonamiento por habitaciones operativo, el sistema no puede demostrar la
ventaja del nivel más alto de la hipótesis y, por tanto, no puede cerrar la
contribución del TFG.

\subsection{Inspiración conceptual: MoMa-LLM}

El paper \textit{Language-Grounded Dynamic Scene Graphs for Interactive Object
Search with Mobile Manipulation} (Honerkamp et al., RSS SemRob~2024) propone
usar un LLM para razonar sobre un grafo de escena estructurado por
habitaciones. Este proyecto adopta esa idea conceptualmente: habitaciones como
entidades explícitas, estado estructurado para el LLM, historial compacto de
búsqueda y selección iterativa de la siguiente habitación.

\textbf{No se utiliza ningún código del repositorio MoMa-LLM}. La implementación
es completamente propia. La diferencia principal es que MoMa-LLM construye el
grafo dinámicamente durante la exploración, mientras que este proyecto usa un
VLMap pre-construido y habitaciones pre-anotadas (HSSD), lo que permite una
evaluación más controlada y reproducible.

% =============================================================================
\section{Estado actual del sistema}
% =============================================================================

\subsection{Resumen ejecutivo}

A día de hoy el sistema dispone de toda la infraestructura, todos los
componentes semánticos y la mayor parte de la maquinaria room-aware. El
pipeline M2 (navegación interactiva sobre HSSD con verificación YOLOE)
funciona de extremo a extremo, y las fases preparatorias del razonamiento por
habitaciones (A,~B,~C de la Sección~\ref{sec:room-aware}) están completas.

Lo que aún no está cerrado es el \textbf{bucle de decisión} que convierte ese
material en una política de búsqueda jerárquica dominante: actualmente el
sistema usa los priors de habitación como información auxiliar, pero la
selección final del candidato sigue ocurriendo de forma esencialmente plana
sobre el heatmap. Esa es la pieza pendiente, y es exactamente la diferencia
decisiva que permitirá demostrar la hipótesis central del TFG.

\subsection{Componentes completados}

\subsubsection*{Infraestructura}

\begin{itemize}
  \item[\checkmark] Docker con CUDA~12.1 + Habitat-Sim~0.3.1 + conda
        (Python~3.9).
  \item[\checkmark] 3 escenas MP3D con VLMaps construidos:
        \textbf{gTV8FGcVJC9\_1}, \textbf{jh4fc5c5qoQ\_1},
        \textbf{mJXqzFtmKg4\_1}.
  \item[\checkmark] Escena HSSD integrada: \textbf{102344280}
        (15~habitaciones, $\sim$447\,m$^2$), VLMap construido e indexado.
  \item[\checkmark] Menú interactivo en Docker (\texttt{menu.sh}) para toda la
        pipeline.
  \item[\checkmark] Submódulo \texttt{third\_party/vlmaps} (rama
        \texttt{tfg-modifications}).
  \item[\checkmark] Repositorio limpio en GitHub (eliminado el binario de
        $571$\,MB del historial con \texttt{git-filter-repo}).
\end{itemize}

\subsubsection*{Pipeline M2 (búsqueda plana funcional)}

\begin{itemize}
  \item[\checkmark] Flujo completo instrucción $\rightarrow$ LLM $\rightarrow$
        selección de goal $\rightarrow$ planificación $\rightarrow$ navegación
        live con visualización.
  \item[\checkmark] Planificador A* \emph{clearance-aware} sobre el navmesh
        de Habitat. La colisión física se delega a la propia navmesh
        (\textit{sliding}); un watchdog de no-progreso aborta si el robot
        deja de avanzar (ver diario del 16~y~20 de abril).
  \item[\checkmark] \texttt{plan\_path\_only()} muestra el camino antes de
        ejecutarlo (línea roja, $800$\,ms de previsualización).
  \item[\checkmark] Multi-goal retry: hasta 8~candidatos por categoría con
        descarte por seguridad y por intentos previos.
\end{itemize}

\subsubsection*{Componentes semánticos}

\begin{itemize}
  \item[\checkmark] \textbf{Mapper LLM}: GPT-4o-mini traduce el objeto a
        categorías VLMap probables y a habitaciones plausibles. Distingue
        queries directas (\texttt{sofa} $\in$ categorías HSSD) e indirectas
        (\texttt{laptop} $\notin$ categorías HSSD). Ejemplo:
        \texttt{laptop} $\rightarrow$
        \{\texttt{categories}: \texttt{[table, counter, shelving]},
        \texttt{rooms}: \texttt{[office, bedroom]}\}.
  \item[\checkmark] \textbf{YOLOE como verificador}: worker persistente,
        protocolo stdout estable, scan~$360^\circ$ asíncrono, visual servoing
        iterativo, caché global por target. Verificación en tres etapas
        (standoff, scan, alternativa). 15~categorías incompatibles
        (\texttt{wall}, \texttt{floor}, \texttt{ceiling}, \ldots) saltan la
        verificación.
  \item[\checkmark] \textbf{Heatmap semántico} con postproceso de calidad por
        componente, filtros de coherencia, área y pared, y fórmula de
        \textit{quality} (Sección~\ref{sec:heatmap-flow}).
\end{itemize}

\subsubsection*{Capa room-aware ya disponible (Fases A, B, C, D)}

\begin{itemize}
  \item[\checkmark] \textbf{Fase A --- capa de habitaciones HSSD}.
        \texttt{SemanticSceneRoomProvider} lee el
        \texttt{semantic\_config.json}, transforma los polígonos de Habitat al
        grid VLMap y proporciona \texttt{get\_room\_at\_cell(row,col)},
        \texttt{get\_room\_centroid(name)} y \texttt{list\_rooms()}.
  \item[\checkmark] \textbf{Fase B --- estado estructurado}. \texttt{RoomState}
        registra para cada habitación: centroide, celdas totales y libres,
        ratio explorado, objetos vistos, candidatos probados/confirmados,
        relevancia para el target actual. \texttt{SearchState} agrega todos
        los \texttt{RoomState}, mantiene el historial de visitas y la
        habitación actual.
  \item[\checkmark] \textbf{Fase C --- priors objeto $\rightarrow$ habitación}.
        Fusión ponderada de tres fuentes (LLM, tabla manual de
        $\sim$35~entradas, evidencia de co-ocurrencia) en
        \texttt{compute\_room\_priors()}. El resultado normalizado se almacena
        en \texttt{RoomState.target\_relevance} y se imprime como traza al
        inicio de cada búsqueda.
  \item[\checkmark] \textbf{Fase D --- selector de habitación}. Antes del
        ranking final de candidatos, el sistema calcula un
        \texttt{room\_score} por habitación combinando prior semántico,
        evidencia del heatmap, coste BFS desde la posición actual y
        penalización por intentos fallidos. La habitación ganadora restringe
        el conjunto de candidatos que pasa al ranking fino.
\end{itemize}

\subsubsection*{Soporte experimental}

\begin{itemize}
  \item[\checkmark] \texttt{place\_objects.py}: define objetos 3D
        personalizados (\texttt{.glb}) vía JSON. Los objetos quedan visibles
        para YOLOE pero ausentes del VLMap, lo que permite evaluar
        directamente el caso de uso \emph{objetos pequeños no representados en
        el mapa semántico}.
\end{itemize}

\subsection{Tabla resumen de estado}

\begin{center}
\begin{tabular}{ll}
\toprule
\textbf{Componente} & \textbf{Estado} \\
\midrule
Infraestructura Docker + Habitat + VLMaps & \textcolor{done}{Completado} \\
Pipeline M2 (planificación + ejecución) & \textcolor{done}{Completado} \\
Mapper LLM (objeto $\rightarrow$ categorías + rooms) & \textcolor{done}{Completado} \\
YOLOE persistente + scan + visual servoing & \textcolor{done}{Completado} \\
Filtros de calidad y heatmap & \textcolor{done}{Completado} \\
Integración HSSD completa & \textcolor{done}{Completado} \\
Inserción de objetos personalizados & \textcolor{done}{Completado} \\
Fase A: capa de habitaciones HSSD & \textcolor{done}{Completado} \\
Fase B: \texttt{RoomState} + \texttt{SearchState} & \textcolor{done}{Completado} \\
Fase C: priors objeto $\rightarrow$ habitación & \textcolor{done}{Completado} \\
\midrule
Fase D: selector de habitación & \textcolor{done}{Completado} \\
Fase E: candidatos intra-habitación & \textcolor{done}{Completado} \\
Fase F: acciones de alto nivel & \textcolor{done}{Completado} \\
Fase G: LLM estratégico con estado & \textcolor{partialcol}{Parcial} \\
Fase H: métricas room-aware & \textcolor{partialcol}{Parcial} \\
\midrule
Apartado extra: baseline vs.\ executor & \textcolor{partialcol}{Parcial} \\
Baselines M1, M3 & \textcolor{partialcol}{Posterior} \\
Evaluación cuantitativa amplia & \textcolor{partialcol}{Posterior} \\
Inserción extensiva de objetos pequeños & \textcolor{partialcol}{Posterior} \\
Extensión ROS\,1 + Toyota HSR & \textcolor{partialcol}{Futura} \\
\bottomrule
\end{tabular}
\end{center}

\subsection{Resumen honesto del estado}

El sistema dispone ya de los tres pilares semánticos (habitación, mueble,
objeto) y de toda la infraestructura para usarlos. Las tres decisiones
jerárquicas básicas están integradas: el selector de habitación (Fase~D)
decide qué habitación merece inspección, la selección intra-habitación
(Fase~E) restringe los candidatos al subconjunto relevante, y la capa de
acciones de alto nivel (Fase~F) expone esas operaciones como un vocabulario
discreto (\texttt{go\_to\_room}, \texttt{explore\_room},
\texttt{inspect\_candidate}, \texttt{verify\_target}, \texttt{done}) con
serialización JSON, listo para ser producido por un LLM estratégico.

Lo que falta es cerrar la cabeza de la política room-aware: conectar un LLM
estratégico que emita esas acciones a partir del estado estructurado
(Fase~G) y construir la batería de métricas que demuestre cuantitativamente
la mejora frente a la línea base M2 (Fase~H).

% =============================================================================
\section{Arquitectura funcional actual}
% =============================================================================

\subsection{Pipeline de búsqueda M2}

El pipeline M2 funciona de extremo a extremo y constituye la base sobre la
que operará la política room-aware. Su flujo es:

\begin{enumerate}
  \item \textbf{Instrucción}. El usuario introduce una query en lenguaje
        natural.
  \item \textbf{Parsing LLM}. \texttt{parse\_object\_goal\_instruction()}
        (GPT-4o-mini) extrae las categorías semánticas implicadas.
  \item \textbf{Heatmap semántico}. Para cada categoría,
        \texttt{compute\_heatmap()} genera un mapa de probabilidad 2D a
        partir de los scores CLIP del VLMap.
  \item \textbf{Postproceso}. \texttt{postprocess\_heatmap()} filtra ruido,
        calcula componentes conexos y asigna un score de calidad a cada uno.
  \item \textbf{Selección de goal}. Se elige el candidato con mayor calidad
        efectiva (incluyendo bonus de proximidad y bonus por habitación
        local) y se le añade un anillo de aproximación seguro alrededor del
        centroide.
  \item \textbf{Planificación}. \texttt{plan\_path\_only()} calcula el camino
        con A* clearance-aware sin ejecutarlo y permite previsualizarlo.
  \item \textbf{Ejecución}. El \textit{path follower} con \textit{lookahead}
        recorre el camino sobre el navmesh. La colisión se delega a Habitat
        (sliding sobre el navmesh) y un watchdog de no-progreso aborta si el
        robot deja de avanzar.
  \item \textbf{Verificación YOLOE}. Tres stages: standoff a $0{,}25$\,m,
        scan~$360^\circ$ asíncrono y componente alternativa del heatmap.
  \item \textbf{Retry}. Si falla, se prueba el siguiente candidato (hasta~8);
        si se agota la categoría, se pasa a la siguiente.
\end{enumerate}

\subsection{Arquitectura por capas de \texttt{tfg-sim}}

La organización actual del sistema en \texttt{tfg-sim} puede resumirse como
una arquitectura de tres capas. La idea clave es que la semántica y la
estrategia viven por encima del backend de simulación, de modo que la futura
migración a ROS/HSR sólo tenga que sustituir la capa inferior sin reescribir la
lógica de decisión.

\begin{figure}[htbp]
\centering
\begin{tikzpicture}[
  font=\small,
  layer/.style={
    draw,
    rounded corners=4pt,
    very thick,
    align=center,
    text width=8.4cm,
    minimum height=1.55cm,
    inner sep=6pt
  },
  side/.style={font=\footnotesize\itshape, gray, anchor=west},
  flow/.style={-{Latex[length=3mm]}, thick, gray!70!black}
]
  \node[layer, fill=blue!10, draw=blue!55!black] (l3) at (0,0)
    {\textbf{Decisión \& estrategia} \\[1pt]
     \footnotesize política room-aware $\cdot$ priors $\cdot$ LLM};
  \node[layer, fill=green!10, draw=green!50!black, below=0.55cm of l3] (l2)
    {\textbf{Percepción semántica} \\[1pt]
     \footnotesize VLMap+CLIP $\cdot$ heatmap $\cdot$ candidatos};
  \node[layer, fill=orange!12, draw=orange!60!black, below=0.55cm of l2] (l1)
    {\textbf{Ejecución encarnada} \\[1pt]
     \footnotesize Habitat-Sim $\cdot$ A* clearance-aware $\cdot$ YOLOE};

  \draw[flow] ([xshift=-1cm]l3.south) -- node[left, font=\footnotesize]
    {decisiones} ([xshift=-1cm]l2.north);
  \draw[flow] ([xshift=-1cm]l2.south) -- node[left, font=\footnotesize]
    {goals 2D} ([xshift=-1cm]l1.north);
  \draw[flow] ([xshift=1cm]l1.north) -- node[right, font=\footnotesize]
    {observaciones} ([xshift=1cm]l2.south);
  \draw[flow] ([xshift=1cm]l2.north) -- node[right, font=\footnotesize]
    {evidencia} ([xshift=1cm]l3.south);
\end{tikzpicture}
\caption{Arquitectura por capas de \texttt{tfg-sim}. La capa inferior es la
única que cambia al migrar a ROS/HSR; la semántica y la estrategia se
preservan.}
\label{fig:tfg-sim-three-layer}
\end{figure}

\subsection{Flujo completo del heatmap}
\label{sec:heatmap-flow}

\subsubsection*{Paso 1 --- Scores CLIP por vóxel}

Cada vóxel del mapa 3D tiene un vector de scores CLIP de dimensión
$|\mathcal{C}|$. Para una query dada, se extrae el score $s_i \in [0,1]$ de
cada vóxel $i$ para la categoría objetivo. Solo se retienen vóxeles con
argmax en esa categoría y $s_i > 0{,}3$.

\subsubsection*{Paso 2 --- Proyección 2D (max-pool por columna)}

Los scores 3D se proyectan sobre el plano del suelo tomando el máximo por
celda $(r, c)$. Resultado: heatmap flotante $gs \times gs$ (por defecto
$gs = 960$).

\subsubsection*{Paso 3 --- Difusión por distancia}

Difusión basada en \texttt{distance\_transform\_edt}: caída lineal $\times
0{,}3$ por celda. Celdas con score $< 0{,}15$ se anulan (franja $\sim$3
celdas $=15$\,cm).

\subsubsection*{Paso 4 --- Suavizado Gaussiano}

Gaussian blur $5\times5$, $\sigma=1{,}0$ para suprimir ruido de alta
frecuencia.

\subsubsection*{Paso 5 --- Umbral relativo y componentes conectados}

Binarización con $t = 0{,}5 \times \max(\text{heatmap suavizado})$.
Componentes conexos (8-conectividad). Para cada componente: área, score
máximo, score medio, suma de scores, bounding box, centroide.

\subsubsection*{Paso 6 --- Scoring interno y filtro relativo}

\[
  \textit{score}(r) = \textit{mean\_val}(r) \cdot \log(1 + \textit{area}(r))
\]
Se descartan componentes con
$\textit{score}(r) < 0{,}2 \cdot \max_j \textit{score}(j)$.

\subsubsection*{Paso 7 --- Fórmula de calidad}

Combinación normalizada de soporte semántico real (\texttt{sum\_val}), tamaño
útil del componente, score CLIP medio y densidad. La formulación enfatiza el
soporte real para evitar que activaciones pequeñas pero densas desplacen a
regiones grandes y representativas (ver diario del 20~de~abril).

\subsubsection*{Paso 8 --- Selección del goal}

A partir del centroide del componente ganador se busca un \emph{anillo de
aproximación seguro}: posiciones a $8$--$25$~celdas del centroide,
navegables y con holgura suficiente. La meta de navegación es la mejor
posición de ese anillo, no el centroide en sí, lo que evita rutas que
terminarían dentro de un mueble.

\subsection{Proveedor de habitaciones HSSD}

La clase \texttt{SemanticSceneRoomProvider} (en \texttt{room\_provider.py})
lee el fichero \texttt{semantics/scenes/<id>.semantic\_config.json} de cada
escena HSSD y expone:

\begin{itemize}
  \item \texttt{build(config\_path, gs, hab\_to\_grid)}: lee
        \texttt{region\_annotations}, convierte cada polígono de Habitat
        world a coordenadas del grid VLMap mediante \texttt{hab\_to\_grid},
        y rasteriza con test punto-en-polígono. Resultado:
        \texttt{\_region\_grid[gs, gs]} (int32, 0~$=$~sin etiquetar).
  \item \texttt{get\_room\_at\_cell(row, col)} $\rightarrow$ nombre de la
        habitación o \texttt{None}.
  \item \texttt{get\_room\_centroid(room\_name)} $\rightarrow$ (row, col) del
        centroide de la habitación.
  \item \texttt{list\_rooms()} $\rightarrow$ lista de nombres de todas las
        habitaciones.
\end{itemize}

Cada región almacena: \texttt{id}, \texttt{name}, \texttt{category},
\texttt{centroid [r, c]}, \texttt{poly\_xz}, \texttt{floor\_height}.

\subsection{Archivos clave del proyecto}

\begin{center}
\begin{tabular}{ll}
\toprule
\textbf{Archivo} & \textbf{Rol} \\
\midrule
\texttt{application/interactive\_object\_nav.py} &
  Pipeline de navegación principal \\
\texttt{vlmaps/robot/habitat\_lang\_robot.py} &
  Robot + planificador + \texttt{plan\_path\_only} \\
\texttt{vlmaps/utils/llm\_utils.py} &
  Mapper LLM: objeto $\rightarrow$ categorías + habitaciones \\
\texttt{vlmaps/utils/yoloe\_utils.py} &
  YoloeSession: worker persistente, sync/async \\
\texttt{vlmaps/utils/\_yoloe\_persistent\_worker.py} &
  Subprocess YOLOE aislado \\
\texttt{vlmaps/utils/room\_provider.py} &
  SemanticSceneRoomProvider (HSSD nativo) \\
\texttt{vlmaps/utils/room\_priors.py} &
  Priors objeto $\rightarrow$ habitación (Fase~C) \\
\texttt{vlmaps/utils/search\_state.py} &
  \texttt{RoomState} y \texttt{SearchState} (Fase~B) \\
\texttt{vlmaps/utils/navigation\_utils.py} &
  A* clearance-aware + \texttt{plan\_clearance\_aware\_astar} \\
\texttt{vlmaps/navigator/navigator.py} &
  Adaptador del planificador clearance-aware \\
\texttt{config/object\_goal\_navigation\_cfg.yaml} &
  Config Hydra principal \\
\texttt{docker/menu.sh} &
  Menú interactivo Docker \\
\bottomrule
\end{tabular}
\end{center}

% =============================================================================
\section{Plan: cierre de la política room-aware}
\label{sec:room-aware}
% =============================================================================

\subsection{Visión general}

El sistema actual ya dispone de los tres niveles semánticos por separado
---\textbf{habitación} (HSSD), \textbf{mueble/zona} (VLMap) y \textbf{objeto}
(YOLOE)--- y de la maquinaria de estado por habitación (Fase~B) y de los
priors objeto~$\rightarrow$ habitación (Fase~C). Lo que aún no existe es el
\emph{módulo de decisión} que use esa información para dirigir la búsqueda
de forma jerárquica.

El plan se divide en cinco fases incrementales (D--H). Cada fase aporta una
capacidad concreta sin romper la anterior, y el conjunto convierte el sistema
de búsqueda plana actual en un buscador room-aware completo y evaluable:

\begin{center}
\begin{tabular}{cp{8cm}l}
\toprule
\textbf{Fase} & \textbf{Aportación} & \textbf{Función en la hipótesis} \\
\midrule
D & Selector de habitación &
  Decide \emph{dónde} buscar \\
E & Selección de candidatos intra-habitación &
  Decide \emph{qué} inspeccionar \\
F & Acciones de alto nivel &
  Estructura la política \\
G & LLM estratégico sobre estado estructurado &
  Razona sobre la búsqueda \\
H & Métricas room-aware &
  Demuestra cuantitativamente la mejora \\
\bottomrule
\end{tabular}
\end{center}

\subsection{Por qué Fase~D es el punto de inflexión del proyecto}
\label{sec:fase-d-inflexion}

Sin Fase~D, el sistema puede tener priors de habitación, anotaciones, estado
estructurado y todo el material semántico necesario, y aun así seguir
comportándose como una búsqueda global: el flujo de control desemboca siempre
en \texttt{select\_best\_candidate()} sobre el heatmap, donde la habitación
es información incidental.

Con Fase~D el flujo se invierte. La primera decisión deja de ser «¿cuál es el
mejor componente del heatmap?» y pasa a ser «¿cuál es la mejor habitación
para buscar este objeto ahora?». Solo cuando esa pregunta tiene respuesta se
baja al nivel de candidatos.

Es el cambio mínimo que convierte el sistema en algo cualitativamente
distinto de la línea base M2 y, por tanto, en algo que el TFG puede defender
como contribución experimental. Las fases E, F y G refinan progresivamente
esa política: E cierra la restricción espacial sobre los candidatos, F la
formaliza como vocabulario de acciones y G permite que un LLM la dirija. La
fase H proporciona la batería de métricas con la que se demuestra el valor
del enfoque. Todas dependen de que exista la decisión por habitación.

\subsection{Fase D --- Selector de habitación}

\textbf{Objetivo}. Que el sistema decida \emph{primero} en qué habitación
buscar y solo después qué candidato inspeccionar.

\textbf{Idea}. En cada ciclo de búsqueda, antes de elegir un componente del
heatmap, se calcula una puntuación por habitación que combina: (i)~el prior
semántico para el objeto buscado (Fase~C), (ii)~la evidencia visual del
heatmap dentro de la habitación, (iii)~lo poco que se ha explorado, (iv)~el
coste de llegar y (v)~la penalización acumulada por intentos fallidos
previos:

\[
  \textit{room\_score}(r) =
    w_1\,\textit{prior}(r)
  + w_2\,\textit{evidence}(r)
  + w_3\,\textit{unexplored}(r)
  - w_4\,\textit{cost}(r)
  - w_5\,\textit{penalty}(r)
\]

Significado de cada término:

\begin{itemize}
  \item $\textit{prior}(r)$: \texttt{RoomState.target\_relevance}
        (Fase~C), normalizado.
  \item $\textit{evidence}(r)$: suma de la \textit{quality} de los
        componentes del heatmap cuyo centroide cae en $r$ (función ya
        disponible: \texttt{compute\_heatmap\_room\_evidence()}).
  \item $\textit{unexplored}(r) = 1 - \texttt{explored\_ratio}(r)$.
  \item $\textit{cost}(r)$: distancia BFS desde la posición actual del robot
        al centroide de $r$, normalizada al diámetro del mapa.
  \item $\textit{penalty}(r)$: penalización por candidatos probados sin éxito
        previamente en $r$, leída de \texttt{RoomState}.
\end{itemize}

Pesos iniciales sugeridos: $w_1=0{,}30$, $w_2=0{,}25$, $w_3=0{,}20$,
$w_4=0{,}15$, $w_5=0{,}10$, ajustables empíricamente.

\textbf{Integración}. Una nueva función
\texttt{select\_best\_room(search\_state, robot\_pos, heatmap\_evidence,
obs\_map)} devuelve el nombre de la habitación ganadora y se invoca al inicio
de cada ciclo de búsqueda, antes de \texttt{select\_best\_candidate()}.

\textbf{Lo que cambia operativamente}. Esta fase es la que convierte el
razonamiento por habitaciones en una decisión real del robot. Hasta ahora los
priors de habitación se calculan, se imprimen y se usan como bonus dentro de
la selección de candidato; con Fase~D pasan a ser la primera decisión del
ciclo y restringen el universo de candidatos que se consideran a continuación.

\textbf{Esfuerzo estimado}: 1--2 sesiones.

\subsection{Fase E --- Selección de candidatos intra-habitación}

\textbf{Objetivo}. Una vez decidida la habitación, ejecutar la búsqueda
detallada \emph{dentro} de ella sin reescribir la pipeline existente.

\textbf{Idea}. La pipeline M2 sigue siendo la herramienta correcta para
inspeccionar candidatos individuales: heatmap~$\rightarrow$ goal~$\rightarrow$
planificación~$\rightarrow$ ejecución~$\rightarrow$ YOLOE. Lo único que
cambia es el \emph{universo} de candidatos: ya no son los componentes
globales del heatmap, sino los que caen dentro de la habitación elegida por
Fase~D.

\begin{enumerate}
  \item Filtrar \texttt{kept\_components} a aquellos cuyo centroide cae
        dentro de la habitación seleccionada
        (\texttt{get\_room\_at\_cell()}).
  \item Ordenar los candidatos filtrados por la \texttt{quality} ya
        implementada, posiblemente con bonus de proximidad.
  \item Ejecutar la pipeline M2 sobre el mejor candidato.
  \item Si no hay candidatos sólidos en la habitación, explorar el frontier
        más cercano dentro de ella (Fase~F).
  \item Si la query es de ubicación conocida (\textit{``ve a la cocina''}),
        navegar directamente a un goal interior seguro de la habitación
        (mecanismo ya existente).
\end{enumerate}

Esto \textbf{no descarta la pipeline M2}: la encapsula dentro de un bucle
exterior dirigido por habitación. La diferencia operativa con respecto al
sistema actual es que el mismo M2 se ejecuta sobre un subconjunto reducido y
relevante de candidatos, en lugar de sobre el ranking global.

\textbf{Esfuerzo estimado}: 1 sesión.

\subsection{Fase F --- Acciones de alto nivel}

\textbf{Objetivo}. Hacer explícito el vocabulario de acciones del robot, de
forma que la política room-aware pueda manipularlo como tal y no como una
maraña de llamadas de bajo nivel.

\begin{center}
\begin{tabular}{lp{9cm}}
\toprule
\textbf{Acción} & \textbf{Descripción} \\
\midrule
\texttt{go\_to\_room(room)} &
  Navegar a un goal interior seguro de la habitación indicada \\
\texttt{explore\_room(room)} &
  Cobertura sistemática de celdas no visitadas dentro de la habitación \\
\texttt{inspect\_candidate(room, cand)} &
  Aproximación a un candidato concreto + verificación YOLOE \\
\texttt{verify\_target} &
  Scan~$360^\circ$ + visual servoing en la posición actual \\
\texttt{done} &
  Objetivo encontrado, fin de la búsqueda \\
\bottomrule
\end{tabular}
\end{center}

\texttt{go\_to\_room} e \texttt{inspect\_candidate} se apoyan en la
navegación existente. \texttt{explore\_room} es nuevo: genera waypoints que
cubran las celdas navegables no visitadas. Esta capa convierte la política
room-aware en una secuencia de decisiones discretas que un LLM (Fase~G) puede
leer y producir de forma natural.

\textbf{Estado.} \textcolor{done}{Implementado} en el módulo
\texttt{vlmaps/policy/} (commit del 20.04). Contiene:

\begin{itemize}
  \item \texttt{policy/actions.py}: \texttt{ActionType} (enum de cinco
        valores), \texttt{Action} (dataclass inmutable validada en
        \texttt{\_\_post\_init\_\_}) y dos helpers
        \texttt{action\_to\_json}\,/\,\texttt{parse\_action\_json} para el
        intercambio con el LLM (Fase~G). Se publica también
        \texttt{ACTION\_SCHEMA\_PROMPT}, fragmento listo para incrustar en el
        \textit{system prompt}.
  \item \texttt{policy/frontier.py}: \texttt{find\_frontier\_in\_room()}
        elige una celda navegable dentro de una habitación, alejada de las
        celdas ya visitadas (parámetros \texttt{visited\_radius\_cells} y
        \texttt{min\_clearance\_cells}). Resuelve el componente nuevo de
        \texttt{explore\_room}.
  \item Ampliación aditiva de \texttt{SearchState}:
        \texttt{action\_log}, \texttt{visited\_cells} y los métodos
        \texttt{record\_action}, \texttt{record\_visited\_cell},
        \texttt{recent\_actions} (sin tocar la lógica existente, conforme al
        criterio \textit{add-only} del repositorio).
  \item Suite de tests \texttt{tests/test\_phase\_f\_actions.py} (15
        casos): validación de cada \texttt{ActionType}, \textit{round-trip}
        JSON, log de acciones, y \textit{frontier picking} con/sin trail
        previa. Toda la batería pasa sin Habitat ni dependencias pesadas.
\end{itemize}

La capa es \emph{additiva}: el bucle interactivo actual sigue funcionando
sin cambios, y la Fase~G podrá enchufar el ejecutor que despache cada
\texttt{Action} sobre las funciones inline ya validadas
(\texttt{navigate\_to\_room\_stage}, \texttt{select\_best\_candidate},
\texttt{scan\_360\_and\_verify}, \texttt{fine\_visual\_center}).

\subsection{Fase G --- LLM estratégico sobre estado estructurado}

\textbf{Objetivo}. Que el LLM decida \emph{estrategia de búsqueda}, no
trayectoria.

El LLM no recibe contexto bruto ni controla la navegación de bajo nivel:
recibe un \emph{estado estructurado y compacto} y devuelve la habitación que
se debe visitar y la acción de alto nivel a ejecutar. La puntuación de
Fase~D queda como \textit{fallback} numérico cuando el LLM no responda con
claridad o como validación cruzada.

Ejemplo de prompt interno:

\begin{verbatim}
Task: find sink

Current room: hallway

Known rooms:
  - bathroom: objects [mirror, toilet, shower], explored 80%,
              1 candidate tried (vanity-zone: YOLOE no confirmation)
  - kitchen:  objects [counter, cabinet, table], explored 60%,
              0 candidates tried
  - living_room: objects [sofa, table, shelf], explored 90%

Unexplored zones:
  - bathroom: near rear corner
  - kitchen: near counter area

Which room should be searched next, and which category or zone
inside it should be prioritized?
\end{verbatim}

El historial enviado al LLM es compacto: las últimas 5~acciones, las
habitaciones exploradas, los candidatos probados y el motivo de fallo
(no~detectado, heatmap débil, navegación fallida). Esto evita inflar el
prompt y mantiene el coste y la latencia bajo control.

El detalle de cómo se integra esa decisión en el sistema sin tocar el
pipeline actual --- entrypoint nuevo, ejecutor de acciones y comparación
de heatmaps --- queda recogido en el
\hyperref[sec:apartado-extra]{apartado extra} previo a la migración a ROS.

\textbf{Estado.} \textcolor{partialcol}{Implementación inicial lista}. Se ha
añadido un módulo nuevo \texttt{vlmaps/policy/strategic\_policy.py} que
construye un \emph{snapshot} estructurado del episodio (habitación actual,
prior/evidence por habitación, shortlist de candidatos, frontiers sugeridos y
últimas acciones), intenta decidir la \emph{siguiente} acción mediante LLM
(\texttt{gpt-4o-mini}) y valida la salida contra el vocabulario de Fase~F. Si
el LLM falla, devuelve JSON inválido o propone una acción inconsistente, la
decisión cae automáticamente al \emph{fallback} heurístico. El
\texttt{interactive\_object\_nav\_executor.py} deja así de preplanificar una
lista fija de acciones y pasa a operar en un bucle de
\texttt{choose\_next\_action()} $\rightarrow$ \texttt{execute\_action()}, con
modo seleccionable \texttt{heuristic}/\texttt{hybrid}/\texttt{llm} desde el
menú del contenedor. Queda pendiente la validación sistemática en runtime y el
ajuste fino de la política, pero la transición arquitectónica ya está
materializada.

\textbf{Esfuerzo restante estimado}: 1--2 sesiones.

\subsection{Fase H --- Métricas que demuestran la hipótesis}
\label{sec:fase-h}

\textbf{Objetivo}. Medir cuantitativamente si el razonamiento por
habitaciones mejora la búsqueda. Esta fase cierra el bucle del TFG: sin ella,
las fases D--G son una propuesta arquitectónica, pero no una contribución
demostrada.

\textbf{Métricas estándar de navegación semántica:}

\begin{itemize}
  \item \textbf{SR} (\textit{Success Rate}): fracción de episodios en los
        que YOLOE confirma el objetivo.
  \item \textbf{SPL} (\textit{Success weighted by Path Length}): SR
        ponderado por la eficiencia del camino respecto al óptimo.
  \item \textbf{Wrong Visits}: número de candidatos visitados sin detección
        positiva antes del éxito.
  \item \textbf{Time to Find}: pasos de simulación hasta la confirmación.
\end{itemize}

\textbf{Métricas específicas de razonamiento por habitaciones} (las que
separan M2 de room-aware):

\begin{itemize}
  \item \textbf{CFR} (\textit{Correct First Room}): ¿fue la primera
        habitación visitada la que contenía el objetivo?
  \item \textbf{CT2R} (\textit{Correct Top-2 Rooms}): ¿estaba la
        habitación correcta entre las dos primeras visitadas?
  \item \textbf{Rooms Before Success}: número de habitaciones visitadas
        antes de encontrar el objeto.
  \item \textbf{Inspections Before Success}: número de candidatos
        inspeccionados antes del éxito.
\end{itemize}

CFR, CT2R, \textit{Rooms Before Success} e \textit{Inspections Before
Success} son las métricas que diferencian directamente un sistema room-aware
de un sistema plano. Si el razonamiento por habitaciones aporta valor, debe
reflejarse en ellas.

\textbf{Esfuerzo estimado}: 1 sesión.

\subsection{Orden de implementación}

\begin{center}
\begin{tabular}{clll}
\toprule
\textbf{Fase} & \textbf{Descripción} & \textbf{Depende de} & \textbf{Esfuerzo} \\
\midrule
A & Capa de habitaciones HSSD & Infraestructura & \textcolor{done}{Hecho} \\
B & RoomState + SearchState & A & \textcolor{done}{Hecho} \\
C & Priors objeto $\rightarrow$ habitación & Mapper LLM & \textcolor{done}{Hecho} \\
D & Selector de habitación & B + C & \textcolor{done}{Hecho} \\
E & Candidatos intra-habitación & D + pipeline M2 & \textcolor{done}{Hecho} \\
F & Acciones de alto nivel & E & \textcolor{done}{Hecho} \\
G & LLM estratégico con estado & B + C + F & \textcolor{partialcol}{Parcial} \\
H & Métricas de evaluación & F & 1 sesión \\
\midrule
  & \textbf{Esfuerzo restante estimado} & & \textbf{2 sesiones} \\
\bottomrule
\end{tabular}
\end{center}

D y E son el cambio operativo de comportamiento: tras estas dos fases el
robot ya busca de forma jerárquica. F y G estructuran y delegan esa
política. H la evalúa y permite contrastarla con la línea base M2. Esa
secuencia es el camino mínimo y suficiente para cerrar la contribución del
TFG.

% =============================================================================
\section{Fase posterior: baselines, evaluación y nuevos objetos}
% =============================================================================

Una vez consolidada la política room-aware (fases D--H), el proyecto entra
en la fase de evaluación amplia y de ampliación del catálogo de objetos
buscados.

\subsection{Baselines M1 y M3}

\begin{itemize}
  \item \textbf{M1 --- Navegación aleatoria}: el robot elige celdas
        navegables al azar como goals. Sirve como cota inferior de
        rendimiento.
  \item \textbf{M3 --- Ranking global por calidad}: candidatos ordenados por
        la \texttt{quality} ya implementada, sin razonamiento por
        habitaciones. Mide si la calidad semántica del heatmap aporta sobre
        la proximidad pura (M2) y, junto con el sistema room-aware, permite
        atribuir la mejora al nivel correcto de la jerarquía.
\end{itemize}

Ambos baselines se implementan como variantes del pipeline existente y se
comparan directamente con el sistema room-aware sobre las mismas métricas
de la Sección~\ref{sec:fase-h}.

\subsection{Evaluación cuantitativa}

Script automático que ejecuta $N$~trials sobre objetos conocidos y produce:

\begin{itemize}
  \item CSV con resultados por episodio (SR, SPL, Wrong Visits, CFR, CT2R,
        Rooms Before Success, etc.).
  \item Gráficas comparativas M1~vs.\ M2~vs.\ M3~vs.\ room-aware.
  \item Tablas para la memoria del TFG.
\end{itemize}

Para HSSD se utilizan las anotaciones nativas del
\texttt{semantic\_config.json} para generar trials con \emph{ground truth}
automatizado, sin etiquetado manual.

\subsection{Objetos pequeños no representados en el VLMap}

Usar \texttt{tools/place\_small\_objects.py} (sucesor de
\texttt{place\_objects.py}, con \emph{deconfliction} de hosts y
\emph{grid offset} de 35\,cm) para insertar objetos de interés que no
aparecen en VLMap. Catálogo oficial tras la criba del 27.04:

\begin{itemize}
  \item \textbf{Cocina (encimera)}: \texttt{bottle}, \texttt{mug},
        \texttt{teapot}, \texttt{toaster}, \texttt{coffee maker}.
  \item \textbf{Multi-room}: \texttt{laptop} (office/bedroom/living),
        \texttt{book} (bedroom/office/living).
  \item Colocados en posiciones conocidas para evaluación con
        \emph{ground truth} a nivel de habitación.
  \item Visibles para YOLOE pero ausentes del VLMap pre-construido.
\end{itemize}

Quedan explícitamente fuera \texttt{kettle} (morfología confundible con
\texttt{teapot}, mismo prior y mismas habitaciones; su presencia
contaminaba la matriz de confusión sin medir nada nuevo) y
\texttt{trash bin} (aparece en demasiadas habitaciones sin prior
dominante, lo que vuelve indistinguible la búsqueda \emph{room-aware} de
una búsqueda plana y mide más al detector que al razonamiento).

Este conjunto de pruebas es el más representativo de la hipótesis del
proyecto: para encontrar una taza, el robot tiene que razonar que vive
típicamente en la cocina, sobre encimeras o mesas, y verificar visualmente
con YOLOE. Un sistema sin los tres niveles falla por construcción en este
escenario.

\paragraph{Próximo paso: \emph{N}~variantes del mapa por escena.}
La métrica actual evalúa una colocación canónica (objeto en su zona más
probable) y mide si el robot la resuelve. La fase siguiente, ya prevista
en el plan, consiste en generar \emph{N}~versiones del mismo mapa con
los mismos objetos colocados en posiciones plausibles distintas (p.\,ej.\
botella en encimera, en mesa de comedor o en mesilla; libro en
escritorio, mesilla o sofá). El robot ejecutará la misma query sobre las
\emph{N}~variantes y se medirá la \emph{degradación graceful} del
sistema cuando el objeto no está en la habitación más probable. Es la
prueba operativa de la hipótesis: un sistema room-aware debe degradar
suavemente al moverse del prior al \emph{long tail}, mientras que un
sistema plano debería degradar abruptamente. La generación automática
de variantes y la agregación cross-variant queda como tarea posterior a
\texttt{post\_placement\_v2}.

% =============================================================================
\section{Apartado extra: comparación baseline vs.\ executor}
\label{sec:apartado-extra}
% =============================================================================

A fecha 22.04 existe ya una \textbf{implementación funcional} del
\emph{executor} y de un \emph{entrypoint} alternativo, pero el apartado sigue
\textbf{parcial}: la infraestructura de comparación ya está montada, pero falta
ejecutarla a escala y consolidar resultados finales. Su objetivo no es
añadir capacidades nuevas, sino \textbf{poder comparar de forma controlada
la política \emph{inline} actual contra la nueva política basada en
acciones tipadas}, una vez ambas están funcionando. Se aborda únicamente
cuando todo el pipeline room-aware (Fases F y G) y las métricas (Fase H)
estén operativos.

\subsection*{Motivación}

\texttt{interactive\_object\_nav.py} acumula $\approx$2500 líneas de
lógica \emph{inline} validada contra el simulador. Reescribir ese fichero
para introducir la nueva política sería un riesgo de regresión enorme y
haría imposible comparar versiones. La estrategia es la opuesta:
\textbf{no se toca el baseline}. Se crea una pista paralela que reutiliza
las mismas primitivas de bajo nivel y solo cambia la orquestación.

\subsection*{Estructura propuesta}

\begin{itemize}
  \item \texttt{application/interactive\_object\_nav.py} --- se queda como
        \textbf{baseline} estable. Mismo arranque, mismo \emph{setup},
        misma UI, mismo bucle \emph{inline} actual.
  \item \texttt{application/interactive\_object\_nav\_executor.py} ---
        \textbf{entrypoint nuevo}. Reutiliza el arranque, el \emph{setup}
        y la UI del baseline, pero su bucle principal pasa por
        \texttt{choose\_next\_action()} $\rightarrow$
        \texttt{execute\_action(...)}. No duplica primitivas. Su primera
        versión ya está implementada y expuesta en el menú interactivo del
        contenedor.
  \item \texttt{vlmaps/policy/executor.py} --- \textbf{módulo nuevo}.
        Implementa \texttt{execute\_action(robot, action, context, ...)} y
        despacha cada \texttt{Action} sobre las primitivas existentes:
        \texttt{go\_to\_room} $\rightarrow$ navegación a \emph{room goal};
        \texttt{explore\_room} $\rightarrow$ \textit{frontier picker} de
        Fase~F; \texttt{inspect\_candidate} $\rightarrow$ selección de
        candidato + \emph{approach goal}; \texttt{verify\_target}
        $\rightarrow$ YOLOE + scan~$360^\circ$ + centrado visual;
        \texttt{done} $\rightarrow$ cierre con \emph{overlays}.
  \item Lo que ya funciona (parser de instrucción, heatmap, selección de
        candidato, navegación a \emph{room goal}, \emph{approach goal},
        YOLOE, scan~$360^\circ$, centrado visual, \textit{overlays}) se
        \textbf{reutiliza tal cual}; el ejecutor solo lo invoca con otro
        orden y otra política de selección.
  \item Más adelante, el LLM se conecta únicamente como productor de
        \texttt{Action} JSON dentro de \texttt{choose\_next\_action()},
        con la heurística de Fase~D como \emph{fallback}. Nada se vuelve
        monolítico.
\end{itemize}

\subsection*{Qué se compara}

Tener los dos entrypoints conviviendo permite ahora un diseño de evaluación
\textbf{2x2} con dos ejes independientes:

\begin{enumerate}
  \item \textbf{Eje de orquestador}: baseline
        (\texttt{interactive\_object\_nav.py}) frente a executor
        (\texttt{interactive\_object\_nav\_executor.py}).
  \item \textbf{Eje de heatmap}: heatmap baseline
        (mapa crudo antes del postproceso) frente a heatmap postprocesado
        (\texttt{postprocess\_heatmap}).
\end{enumerate}

Combinando ambos ejes aparecen cuatro métodos comparables sobre el mismo JSONL:
\begin{center}
\begin{tabular}{ll}
\toprule
\textbf{Clave} & \textbf{Configuración} \\
\midrule
\texttt{Ob\_Hb} & orquestador baseline + heatmap baseline \\
\texttt{Oe\_Hb} & executor + heatmap baseline \\
\texttt{Ob\_Hp} & orquestador baseline + heatmap postprocesado \\
\texttt{Oe\_Hp} & executor + heatmap postprocesado \\
\bottomrule
\end{tabular}
\end{center}

\subsection*{Criterio de cierre}

El apartado se da por cerrado cuando:
(i)~el baseline sigue funcionando sin regresiones tras introducir el
ejecutor;
(ii)~existe un informe lado a lado de los dos entrypoints sobre el mismo
conjunto de \emph{queries};
(iii)~se publican las imágenes del heatmap baseline y del heatmap
procesado para al menos una escena HSSD representativa.

A partir de ahí, la migración a ROS\,1\,/\,Toyota HSR (sección
siguiente) puede acometerse sustituyendo \emph{solo} las capas de
percepción, navegación y ejecución, sin tocar la capa de decisión, que
ya está validada en simulación.

\textbf{Estado actual} (22.04):
\textcolor{partialcol}{las baterías 2x2 y de heatmaps están integradas; falta lanzarlas
sobre escenas representativas y consolidar las métricas}. La infraestructura
queda dividida en dos pistas independientes que se pueden ejecutar por separado
o en cadena:
\begin{itemize}
  \item \textbf{Generador de \emph{queries} compartido}.
        \texttt{tools/build\_eval\_queries.py} produce \emph{batches} en
        formato JSONL bajo \texttt{tools/eval\_queries/\{scene\}.jsonl}
        (50~queries por escena, 4 tipos: \texttt{object}, \texttt{room},
        \texttt{room\_object}, \texttt{compound}). Cada entrada lleva
        \texttt{id}, \texttt{target\_label}, \texttt{expected\_rooms} y
        \texttt{expected\_room\_polygons} (referencia exacta a polígonos
        del \texttt{semantic\_config.json} para \emph{ground truth}
        espacial). Ese mismo fichero alimenta los dos comparadores.
  \item \textbf{Runner online 2x2}.
        \texttt{tools/run\_nav\_eval.py} acepta ahora, además del
        \emph{entrypoint}, un \texttt{heatmap\_mode}
        (\texttt{baseline} o \texttt{postprocessed}) que se propaga como
        variable de entorno a los dos navegadores. Cada ejecución produce
        \texttt{raw\_log.txt}, \texttt{setup.log}, un directorio
        \texttt{segments/} con un \texttt{.log} por query y un
        \texttt{manifest.json} enriquecido con \texttt{scene\_id},
        \texttt{scene\_name}, \texttt{heatmap\_mode} y toda la metadata
        del JSONL. Encima de esta pieza, \texttt{tools/run\_full\_2x2\_eval.py}
        orquesta automáticamente las cuatro combinaciones
        \texttt{Ob\_Hb}, \texttt{Oe\_Hb}, \texttt{Ob\_Hp} y \texttt{Oe\_Hp},
        copiando además el JSONL bajo
        \texttt{/workspace/results/eval\_runs/\{timestamp\}/queries/}.
  \item \textbf{Agregador de métricas 2x2}.
        \texttt{tools/aggregate\_full\_2x2\_eval.py} consume los cuatro
        manifiestos y produce dos vistas: (i)~\texttt{compare\_full.csv}
        y \texttt{compare\_full.md} con una fila por query y columnas para
        las cuatro combinaciones, y (ii)~\texttt{aggregate\_metrics.csv}
        y \texttt{aggregate\_metrics.md} con métricas agregadas por
        método, por \texttt{query\_type} y por \texttt{tag}. Las métricas
        actualmente calculadas de forma robusta son:
        SR, Object~SR, Room~SR, CFR, CT2R, Rooms~Before~Success,
        Inspections~Before~Success, Wrong~Visits, Mean~Pose~Updates,
        Early~Stop~Rate, Final~Room~Accuracy, Preview~No~Action~Rate,
        Path~Stopped~Early~Rate, Mean~Executor~Actions y
        Mean~Room~Transitions. Para ello, tanto baseline como executor
        emiten al final de cada query un marcador estructurado
        \texttt{[eval-summary] \{...\}} con historial de habitaciones,
        candidatos intentados/confirmados y rastro de acciones.
  \item \textbf{Comparador offline de heatmaps}.
        \texttt{tools/eval\_heatmap\_postprocess.py} sigue calculando,
        para cada query, el \emph{heatmap} crudo (CLIP + proyección 2D +
        decay) y el \emph{heatmap} postprocesado, pero ahora añade además
        agregación por \texttt{query\_type} y por \texttt{tag} en
        \texttt{aggregate\_by\_slice.csv}. El \texttt{summary.md}
        incluye tanto la tabla global como tablas resumidas por tipo de
        query y por etiqueta (\texttt{multi\_instance},
        \texttt{ambiguous\_room}, etc.), mientras que \texttt{per\_query.csv}
        conserva la traza completa query-a-query.
  \item \textbf{Submenú de evaluación simplificado}.
        El menú interactivo del contenedor concentra ahora la evaluación
        bajo una única opción \texttt{t) Testing / evaluation submenu} con
        exactamente tres entradas visibles:
        \texttt{Generate test set},
        \texttt{Compare full 2x2 pipeline} y
        \texttt{Heatmap-only offline analysis}. El usuario ya no tiene que
        encadenar scripts manualmente: basta elegir escena(s) y, si
        quiere, un JSONL alternativo; el \emph{tooling} crea la estructura
        estándar de resultados bajo
        \texttt{/workspace/results/eval\_runs/\{timestamp\}/}.
\end{itemize}
Como ambas pistas comparten \texttt{id}, \texttt{scene\_id},
\texttt{query\_type}, \texttt{expected\_rooms} y \texttt{tags}, los resultados
se pueden cruzar a posteriori para responder preguntas del tipo \emph{«¿las
queries que el executor falla son justo las que ya tenían un
\texttt{Hit@1} bajo en el heatmap raw, o hay fallos genuinos de
orquestación?»}. El trabajo restante ya no es infraestructura, sino lanzar las
baterías sobre escenas HSSD representativas, revisar salidas y comentar
resultados.

\subsection*{Protocolo de evaluación finalmente adoptado}
\label{sec:eval-protocol}

La batería 2$\times$2 inicial se descartó por ser muy costosa de lanzar y
porque mezclaba dos preguntas distintas. Se reemplazó por un \textbf{protocolo
secuencial en dos etapas} que aísla cada pregunta y reduce el tiempo de
cómputo manteniendo la potestad de responder ambas:

\begin{enumerate}
  \item \textbf{Etapa~1 --- selección del heatmap} \emph{(offline, sin
        navegación)}. Sobre 500~queries por escena se compara el heatmap
        \texttt{raw} (CLIP $\rightarrow$ proyección 2D $\rightarrow$
        \emph{decay}) con el heatmap \texttt{postprocessed}
        (\texttt{postprocess\_heatmap()}). El criterio de victoria es
        agregado: \emph{Hit@1}, \emph{Hit@5}, \emph{IoU top-mass-50},
        masa proyectada dentro de la habitación esperada y proporción de
        masa fuera. La salida fija el modo a usar en la siguiente etapa.
  \item \textbf{Etapa~2 --- selección del orquestador} \emph{(online, con
        Habitat-Sim)}. Sobre 100~queries por escena, con el heatmap fijado
        al ganador de la Etapa~1, se compara el orquestador \emph{baseline}
        (\texttt{Ob\_Hp}) con el \emph{executor} de acciones tipadas
        (\texttt{Oe\_Hp}). Aquí las métricas son de navegación: SR, CFR,
        Wrong Visits, Pose Updates y Early Stop Rate.
\end{enumerate}

Esta separación permite que la decisión sobre el heatmap se tome con un
volumen mucho mayor de queries (500 vs.\ 100) y a coste cero de simulación,
mientras que la simulación se reserva para el contraste de orquestador, que
es lo que realmente exige ejecutarse en Habitat.

\subsubsection*{Etapa~1: heatmap \texttt{raw} vs.\ \texttt{postprocessed}}

Lote de referencia: \texttt{results/heatmap\_phase/20260424\_153612/}
(500 queries $\times$ 3 escenas HSSD). La Tabla~\ref{tab:heatmap-stage}
recoge los valores por escena y la Figura~\ref{fig:heatmap-stage-quality}
muestra las cuatro métricas clave por escena.

\begin{table}[htbp]
\centering
\small
\setlength{\tabcolsep}{4pt}
\begin{tabular}{lrrrrrrrr}
\toprule
\multirow{2}{*}{\textbf{Escena}} & \multirow{2}{*}{\textbf{n}} &
\multicolumn{2}{c}{\textbf{Hit@1}} &
\multicolumn{2}{c}{\textbf{Mass in expected}} &
\multicolumn{2}{c}{\textbf{Wrong-room mass}} \\
\cmidrule(lr){3-4} \cmidrule(lr){5-6} \cmidrule(lr){7-8}
 & & raw & postp. & raw & postp. & raw & postp. \\
\midrule
102344193\_0          & 500 & 0.430 & \textbf{0.430} & 0.374 & \textbf{0.400} & \textbf{0.564} & 0.589 \\
102344280\_0          & 500 & 0.316 & \textbf{0.414} & 0.281 & \textbf{0.310} & 0.577 & \textbf{0.551} \\
108736884\_177263634\_4 & 500 & 0.326 & \textbf{0.352} & 0.276 & \textbf{0.298} & 0.611 & \textbf{0.604} \\
\bottomrule
\end{tabular}
\caption{Etapa~1: comparativa offline del heatmap \texttt{raw} frente al
\texttt{postprocessed} sobre 500 queries por escena. \emph{Hit@1} mide si la
celda de máxima masa cae dentro de la habitación esperada; \emph{mass in
expected} es la fracción de masa total proyectada dentro de las habitaciones
correctas; \emph{wrong-room mass} es la masa proyectada fuera (cuanto menor,
mejor). Negrita = mejor entre los dos modos.}
\label{tab:heatmap-stage}
\end{table}

\begin{figure}[htbp]
\centering
\includegraphics[width=0.49\linewidth]{vlmap-semantic-object-search-tfg/results/heatmap_phase/20260424_153612/plots/bar_hitat_1.png}\hfill
\includegraphics[width=0.49\linewidth]{vlmap-semantic-object-search-tfg/results/heatmap_phase/20260424_153612/plots/bar_mass_in_expected_ratio.png}\\[6pt]
\includegraphics[width=0.49\linewidth]{vlmap-semantic-object-search-tfg/results/heatmap_phase/20260424_153612/plots/bar_iou_topmass50.png}\hfill
\includegraphics[width=0.49\linewidth]{vlmap-semantic-object-search-tfg/results/heatmap_phase/20260424_153612/plots/bar_wrong_room_mass_ratio.png}
\caption{Etapa~1: calidad semántica del heatmap por escena. \emph{Hit@1}
(arriba-izq.), \emph{mass in expected room} (arriba-der.), \emph{IoU top-mass-50}
(abajo-izq.) y \emph{wrong-room mass} (abajo-der.). Cada barra compara el modo
\texttt{raw} con el modo \texttt{postprocessed}; el postproceso gana en
\emph{Hit@1} y \emph{mass in expected} en las tres escenas.}
\label{fig:heatmap-stage-quality}
\end{figure}

\paragraph{Análisis Etapa~1.} El postprocesado mantiene o mejora \emph{Hit@1}
y \emph{mass in expected} en las tres escenas, mejora \emph{wrong-room mass}
en dos de tres, y empata o cae ligeramente en \emph{IoU top-mass-50} (el
postproceso compacta los componentes y reduce la masa total, lo que penaliza
una métrica de solapamiento de masa). En agregado el efecto neto es positivo,
con la mejora más clara en la escena 102344280\_0 (Hit@1 0.32~$\rightarrow$~0.41,
$+9.8$~puntos). \textbf{Decisión}: se fija \texttt{heatmap\_mode = postprocessed}
para la Etapa~2.

\subsubsection*{Etapa~2: orquestador \texttt{Ob\_Hp} vs.\ \texttt{Oe\_Hp}}

Lote de referencia: \texttt{results/orchestrator\_phase/20260424\_154155/}
(100 queries $\times$ 3 escenas, heatmap fijado a \texttt{postprocessed}).
\texttt{Ob\_Hp} es el orquestador \emph{inline} actual; \texttt{Oe\_Hp} es el
nuevo executor que pasa por \texttt{choose\_next\_action() $\rightarrow$
execute\_action()} sobre las mismas primitivas. Las
Tablas~\ref{tab:orch-stage-cross} y~\ref{tab:orch-stage-perscene} resumen los
resultados; las Figuras \ref{fig:orch-sr-cfr} y \ref{fig:orch-costs} muestran
las métricas en formato gráfico por escena.

\begin{table}[htbp]
\centering
\small
\begin{tabular}{lrrrrr}
\toprule
\textbf{Método} & \textbf{SR} & \textbf{CFR} & \textbf{CT2R} & \textbf{Wrong Visits} & \textbf{Pose Updates} \\
\midrule
\texttt{Ob\_Hp} (baseline) & \textbf{0.667 $\pm$ 0.18} & \textbf{0.137 $\pm$ 0.07} & \textbf{0.290 $\pm$ 0.07} & \textbf{0.16 $\pm$ 0.12} & \textbf{382.4 $\pm$ 19.8} \\
\texttt{Oe\_Hp} (executor) & 0.617 $\pm$ 0.14 & 0.177 $\pm$ 0.07 & 0.327 $\pm$ 0.07 & 1.18 $\pm$ 0.44 & 431.6 $\pm$ 35.9 \\
\midrule
$\Delta$ (Oe$-$Ob) & $-0.050$ & $+0.040$ & $+0.037$ & $+1.02$ & $+49.1$ \\
\bottomrule
\end{tabular}
\caption{Etapa~2: agregado cruzado de las 3 escenas (300 queries). SR =
Success Rate, CFR = Completion Failure Rate, CT2R = Cost to Result. Negrita
= mejor de los dos métodos. El baseline gana en las cinco columnas.}
\label{tab:orch-stage-cross}
\end{table}

\begin{table}[htbp]
\centering
\small
\setlength{\tabcolsep}{5pt}
\begin{tabular}{llrrrrr}
\toprule
\textbf{Escena} & \textbf{Método} & \textbf{SR} & \textbf{CFR} & \textbf{Wrong Visits} & \textbf{Pose Updates} & \textbf{Early Stop} \\
\midrule
\multirow{2}{*}{102344193\_0}
 & \texttt{Ob\_Hp} & \textbf{0.92} & 0.24 & \textbf{0.00} & \textbf{404.0} & \textbf{0.21} \\
 & \texttt{Oe\_Hp} & 0.81          & \textbf{0.27} & 0.56          & 381.2          & 0.04 \\
\midrule
\multirow{2}{*}{102344280\_0}
 & \texttt{Ob\_Hp} & 0.50 & \textbf{0.09} & \textbf{0.28} & \textbf{387.2} & \textbf{0.02} \\
 & \texttt{Oe\_Hp} & 0.50 & 0.10          & 1.55          & 451.4          & 0.01 \\
\midrule
\multirow{2}{*}{108736884\_177263634\_4}
 & \texttt{Ob\_Hp} & \textbf{0.58} & \textbf{0.08} & \textbf{0.20} & \textbf{356.2} & \textbf{0.07} \\
 & \texttt{Oe\_Hp} & 0.54          & 0.16          & 1.43          & 462.2          & 0.03 \\
\bottomrule
\end{tabular}
\caption{Etapa~2 desglosada por escena. Negrita = mejor de los dos métodos en
esa escena. El baseline gana o empata en SR en las tres escenas, y siempre
gana en \emph{Wrong Visits} y \emph{Pose Updates}.}
\label{tab:orch-stage-perscene}
\end{table}

\begin{figure}[htbp]
\centering
\includegraphics[width=0.49\linewidth]{vlmap-semantic-object-search-tfg/results/orchestrator_phase/20260424_154155/plots/sr.png}\hfill
\includegraphics[width=0.49\linewidth]{vlmap-semantic-object-search-tfg/results/orchestrator_phase/20260424_154155/plots/cfr.png}
\caption{Etapa~2: Success Rate (izq.) y Completion Failure Rate (der.) por
escena para \texttt{Ob\_Hp} y \texttt{Oe\_Hp}.}
\label{fig:orch-sr-cfr}
\end{figure}

\begin{figure}[htbp]
\centering
\includegraphics[width=0.49\linewidth]{vlmap-semantic-object-search-tfg/results/orchestrator_phase/20260424_154155/plots/wrong_visits.png}\hfill
\includegraphics[width=0.49\linewidth]{vlmap-semantic-object-search-tfg/results/orchestrator_phase/20260424_154155/plots/mean_pose_updates.png}
\caption{Etapa~2: \emph{Wrong Visits} (izq.) y \emph{Mean Pose Updates}
(der.) por escena. La diferencia más visible no está en el éxito sino en el
coste: el executor inspecciona muchos más candidatos incorrectos antes de
cerrar la query.}
\label{fig:orch-costs}
\end{figure}

\paragraph{Análisis Etapa~2.} El baseline \texttt{Ob\_Hp} gana en SR cruzado
($0.667$ vs.\ $0.617$, $\Delta=-5$~puntos) y en CFR ($0.137$ vs.\ $0.177$). La
diferencia más relevante no es el SR sino el coste: el executor genera
$\sim$7$\times$ más \emph{Wrong Visits} ($1.18$ vs.\ $0.16$) y un 13\,\%
más de \emph{Pose Updates} ($432$ vs.\ $382$). Por escena, el baseline gana
limpiamente en sc1 ($0.92$ vs.\ $0.81$) y sc3 ($0.58$ vs.\ $0.54$), y empata
en sc2 ($0.50$/$0.50$). Por tipo de query, el executor sí mejora ligeramente
en \emph{room\_object} en sc2 ($0.43$ vs.\ $0.37$), donde la decisión tipada
explícita ayuda a aprovechar el \emph{room hint}; en \emph{object} puro
pierde en las tres escenas. \textbf{Decisión}: el executor se mantiene en el
repositorio para iterar sobre el \emph{action vocabulary}, pero el orquestador
de referencia para el resto del proyecto sigue siendo \texttt{Ob\_Hp}, con la
salvedad de explorar el \emph{action vocabulary} para tareas explícitamente
\emph{room-guided}.

\paragraph{Lo que estas dos etapas dejan asentado.}
(i)~El postproceso del heatmap es la elección por defecto; (ii)~el cuello de
botella ya no es el modo de heatmap sino la calidad de la política de
inspección dentro de la habitación; (iii)~el executor sólo recupera al
baseline en escenarios con \emph{room hint} explícito, lo que indica que el
camino para que valga la pena pasa por enriquecer el \emph{action vocabulary}
y conectarle el LLM estratégico (Fase~G), no por reescribir más primitivas.

\subsection*{Calibración del umbral de confianza de YOLOE}
\label{sec:yoloe-thresh}

El verificador visual YOLOE actúa como filtro final del pipeline: cuando el
agente llega a un candidato propuesto por el heatmap, decide si la detección
supera un umbral de confianza ($\tau$) para confirmar el objetivo. Un $\tau$
demasiado bajo aceptaría detecciones débiles (riesgo de falsos positivos);
uno demasiado alto rechazaría detecciones reales (falsos negativos), forzando
al agente a seguir explorando.

Para fijar este parámetro de forma justificada se realizó un barrido
controlado sobre $\tau \in \{0.30, 0.40, 0.50, 0.60\}$ en una única configuración
del 2$\times$2 (executor + heatmap postprocesado, $Oe\_Hp$) y una única escena
($102344193\_0$, 50 queries de muebles y categorías VLMap). Como las decisiones
de navegación, selección de habitación y ranking de candidatos no dependen de
$\tau$, basta con barrer un método para aislar el efecto del verificador; los
otros tres métodos del 2$\times$2 heredarán el valor seleccionado.

\begin{table}[h]
\centering
\small
\begin{tabular}{lccccc}
\toprule
$\tau$ & SR & Obj.\ SR & Found rate & Wrong visits & CFR \\
\midrule
\textbf{0.30} & \textbf{0.82} & \textbf{0.76} & \textbf{0.62} & \textbf{0.84} & 0.18 \\
0.40 & 0.76 & 0.68 & 0.56 & 1.26 & 0.18 \\
0.50 & 0.78 & 0.71 & 0.56 & 1.42 & 0.18 \\
0.60 & 0.66 & 0.55 & 0.44 & 2.63 & 0.18 \\
\bottomrule
\end{tabular}
\caption{Métricas end-to-end para los cuatro umbrales evaluados sobre $Oe\_Hp$
en $102344193\_0$ (50 queries).}
\label{tab:yoloe-thresh}
\end{table}

\begin{figure}[h]
\centering
\begin{tikzpicture}
\begin{axis}[
    width=0.85\linewidth, height=6cm,
    xlabel={Umbral de confianza $\tau$},
    ylabel={Tasa},
    xmin=0.28, xmax=0.62,
    ymin=0.0, ymax=1.0,
    xtick={0.30,0.40,0.50,0.60},
    legend style={at={(1.02,1)}, anchor=north west, font=\small},
    grid=both, grid style={gray!20},
    cycle list name=color,
    mark options={solid},
]
\addplot+[mark=*]   coordinates {(0.30,0.82) (0.40,0.76) (0.50,0.78) (0.60,0.66)};
\addplot+[mark=square*] coordinates {(0.30,0.76) (0.40,0.68) (0.50,0.71) (0.60,0.55)};
\addplot+[mark=triangle*] coordinates {(0.30,0.62) (0.40,0.56) (0.50,0.56) (0.60,0.44)};
\legend{SR, Object SR, Found rate}
\end{axis}
\end{tikzpicture}
\caption{Evolución de SR, Object SR y tasa de detección al aumentar $\tau$.
La degradación es monótona a partir de $\tau=0.40$.}
\label{fig:yoloe-thresh-rates}
\end{figure}

\begin{figure}[h]
\centering
\begin{tikzpicture}
\begin{axis}[
    width=0.85\linewidth, height=5cm,
    xlabel={Umbral de confianza $\tau$},
    ylabel={Visitas erróneas (media)},
    xmin=0.28, xmax=0.62,
    ymin=0.5, ymax=3.0,
    xtick={0.30,0.40,0.50,0.60},
    grid=both, grid style={gray!20},
    mark options={solid},
]
\addplot+[mark=*, color=red] coordinates {(0.30,0.84) (0.40,1.26) (0.50,1.42) (0.60,2.63)};
\end{axis}
\end{tikzpicture}
\caption{Visitas erróneas medias por query (candidatos inspeccionados pero
no confirmados) en función de $\tau$. El crecimiento monótono indica que el
verificador rechaza cada vez más detecciones reales.}
\label{fig:yoloe-thresh-wrong}
\end{figure}

\paragraph{Análisis.}
La métrica CFR (\emph{Correct First Room}) se mantiene constante en $0.18$ para
los cuatro umbrales, lo que confirma que el experimento aísla correctamente el
efecto del verificador: la navegación y la selección de habitación no se ven
alteradas. Las tres métricas que sí dependen del verificador (SR, Object SR y
\emph{found rate}, Figura~\ref{fig:yoloe-thresh-rates}) muestran su valor
máximo en $\tau=0.30$ y caen al aumentar el umbral. La tasa de visitas erróneas
(Figura~\ref{fig:yoloe-thresh-wrong}) crece de forma monótona, llegando a
triplicarse entre $\tau=0.30$ ($0.84$) y $\tau=0.60$ ($2.63$): el agente
inspecciona más muebles candidatos antes de aceptar uno como verificado.

\paragraph{Decisión.}
Se selecciona $\tau=0.30$ por maximizar SR y Object SR y minimizar las visitas
erróneas. Este valor se utiliza en todos los métodos del 2$\times$2 oficial.

\paragraph{Limitaciones.}
La medida directa de la tasa de falsos positivos del verificador requeriría
\emph{queries negativas} (objetivos ausentes en la escena, donde cualquier
\texttt{found=True} sería FP puro) o \emph{ground-truth} a nivel de instancia
para comparar la posición confirmada con la posición real del objetivo.
Ninguno de los dos esquemas estaba disponible en este trabajo y se considera
una línea de mejora del banco de evaluación.

% =============================================================================
\section{Extensión futura: ROS\,1 Noetic + Toyota HSR}
% =============================================================================

El proyecto se desarrolla y valida primero en simulación, con HSSD como
entorno principal y con foco en la capacidad de razonamiento por habitaciones
y búsqueda semántica de objetos. Una vez consolidada la capa de decisión de
alto nivel y evaluado el sistema sobre tareas representativas, se contempla
como línea de extensión la migración a ROS\,1 Noetic y su integración con el
robot Toyota HSR.

La arquitectura de razonamiento por habitaciones (Sección~\ref{sec:room-aware})
está diseñada para ser independiente de la plataforma de ejecución: las
capas de decisión (selector de habitación, LLM estratégico, estado de
búsqueda) no dependen de Habitat-Sim. La migración sustituiría únicamente
las capas de percepción, navegación y ejecución por sus equivalentes en
robot real:

\begin{itemize}
  \item \texttt{vlmap\_server} (nodo ROS): carga el VLMap y expone un
        servicio \texttt{/vlmap/query} para consultas semánticas.
  \item \texttt{vlmap\_nav} (nodo ROS): recibe instrucciones en lenguaje
        natural, ejecuta el razonamiento por habitaciones y envía goals a
        \texttt{move\_base}.
  \item Visualización en RViz: heatmap como \texttt{OccupancyGrid}, camino
        como \texttt{nav\_msgs/Path}, estado de habitaciones como markers.
  \item Cámara RGBD del HSR como entrada para YOLOE (sustituyendo la cámara
        simulada de Habitat-Sim).
\end{itemize}

Esta extensión no forma parte del roadmap inmediato. Su viabilidad depende
de que la capa de razonamiento funcione correctamente en simulación y de la
disponibilidad del robot HSR en el laboratorio V4R del TU Wien.

\subsection*{Esquema futuro de doble contenedor}

La evolución prevista del proyecto no consiste en ``mover todo a ROS'', sino en
mantener dos contenedores con responsabilidades nítidas. El contenedor
\texttt{tfg-sim} seguirá siendo la referencia para HSSD/Habitat, datasets,
evaluación y razonamiento semántico; \texttt{tfg-ros} absorberá ROS\,1 Noetic,
Gazebo, RViz y la interfaz hacia HSR-sim o el robot físico.

\begin{figure}[htbp]
\centering
\begin{tikzpicture}[
  font=\small,
  container/.style={
    draw,
    rounded corners=5pt,
    very thick,
    minimum width=6.4cm,
    minimum height=6.0cm,
    inner sep=4pt
  },
  module/.style={
    draw,
    rounded corners=2pt,
    fill=white,
    align=center,
    text width=5.4cm,
    minimum height=0.65cm,
    font=\footnotesize
  },
  ctitle/.style={font=\bfseries\large, anchor=north},
  flow/.style={-{Latex[length=3mm]}, thick, gray!70!black}
]
  % --- tfg-sim ---
  \node[container, fill=blue!8, draw=blue!55!black] (sim) at (-3.6,0) {};
  \node[ctitle] at (sim.north) [yshift=-0.25cm] {\texttt{tfg-sim}};
  \node[module, draw=blue!50!black] (s1) at (-3.6, 1.85) {VLMap + CLIP};
  \node[module, draw=blue!50!black, below=0.20cm of s1] (s2) {Heatmap (postproc)};
  \node[module, draw=blue!50!black, below=0.20cm of s2] (s3) {Room provider HSSD};
  \node[module, draw=blue!50!black, below=0.20cm of s3] (s4) {Política room-aware};
  \node[module, draw=blue!50!black, below=0.20cm of s4] (s5) {Eval \& métricas};

  % --- tfg-ros ---
  \node[container, fill=green!8, draw=green!50!black] (ros) at (3.6,0) {};
  \node[ctitle] at (ros.north) [yshift=-0.25cm] {\texttt{tfg-ros}};
  \node[module, draw=green!50!black] (r1) at (3.6, 1.85) {ROS\,1 Noetic};
  \node[module, draw=green!50!black, below=0.20cm of r1] (r2) {Gazebo + RViz};
  \node[module, draw=green!50!black, below=0.20cm of r2] (r3) {\texttt{move\_base} + costmaps};
  \node[module, draw=green!50!black, below=0.20cm of r3] (r4) {HSR (sim/real)};
  \node[module, draw=green!50!black, below=0.20cm of r4] (r5) {Bridge sensores};

  % --- bridge between containers ---
  \draw[flow] (sim.east|-s2) -- node[above, font=\footnotesize]
    {goals semánticos} (ros.west|-r3);
  \draw[flow] (ros.west|-r5) -- node[below, font=\footnotesize]
    {odom $\cdot$ RGBD} (sim.east|-s4);
\end{tikzpicture}
\caption{Arquitectura objetivo de doble contenedor. \texttt{tfg-sim} conserva
la pila semántica completa; \texttt{tfg-ros} se encarga de la ejecución sobre
HSR (simulado o físico). El acoplamiento se reduce a goals 2D y observaciones.}
\label{fig:dual-container-architecture}
\end{figure}

% =============================================================================
\section{Diario de desarrollo}
% =============================================================================

\subsection*{31 de marzo de 2026 --- Correcciones del pipeline HSSD}

\begin{enumerate}
  \item \textbf{Mapa de obstáculos invertido}: modelo de dos bandas (suelo $h<0{,}15$\,m /
        obstáculo $0{,}2 \leq h < 1{,}5$\,m). Una celda navegable $= 1$ si tiene suelo Y
        no tiene obstáculo.
  \item \textbf{Heatmap ruidoso}: argmax binario reemplazado por score CLIP continuo con
        umbral 0.3 + max-pooling + difusión solo desde celdas con score alto.
  \item \textbf{Robot atraviesa paredes en HSSD}: \texttt{sim.recompute\_navmesh()} en
        \texttt{\_setup\_sim()} cuando \texttt{dataset\_type == "hssd"}.
  \item \textbf{Integración YOLOE}: \texttt{yoloe\_utils.py} + \texttt{\_yoloe\_worker.py},
        superposición sobre ventana 1st person.
  \item \textbf{Camino visible}: \texttt{path\_cells} pasado a \texttt{show\_map()},
        color rojo.
  \item \textbf{Cámara primera persona}: inclinación $-\pi/8$\,rad ($\approx 22{,}5°$).
\end{enumerate}

\subsection*{2 de abril de 2026 (primera iteración)}

\begin{enumerate}
  \item \textbf{Postproceso de heatmaps}: suavizado Gaussiano $\rightarrow$ umbral
        relativo $\rightarrow$ metadatos por componente (área, max, media, bbox, centroide).
  \item \textbf{Exploración 360°}: pasos de 20°, detención al detectar.
  \item \textbf{Ruta alternativa}: componente con mayor score a $>$20 celdas.
  \item \textbf{Movimiento fluido}: sin \texttt{set\_agent\_state()} visible.
\end{enumerate}

\subsection*{2 de abril de 2026 (segunda iteración)}

\begin{enumerate}
  \item \textbf{Worker YOLOE persistente}: modelo cargado una sola vez; protocolo
        stdin/stdout; inferencias $\sim$100--500\,ms.
  \item \textbf{YoloeSession}: modos síncrono y asíncrono (hilo bg).
  \item \textbf{Scan 360° fluido}: pasos nativos de 5°, YOLOE en hilo bg.
  \item \textbf{Cámara $-30°$}: \texttt{habitat\_utils.py}.
  \item \textbf{Standoff 0,25\,m}: acercamiento máximo al objeto.
\end{enumerate}

\subsection*{13 de abril de 2026}

\begin{enumerate}
  \item \textbf{Bug protocolo stdout/stderr YOLOE}: \texttt{\_ensure\_mobileclip()}
        imprimía en stdout interceptando el \texttt{"READY"}; corregido a
        \texttt{sys.stderr}. \texttt{start()} ahora lee hasta recibir \texttt{"READY"}.
  \item \textbf{mobileclip\_blt.ts}: búsqueda en múltiples rutas candidatas;
        eliminación automática de versión corrupta (${<}100$\,MB).
  \item \textbf{Captura stderr del worker}: crashes visibles en consola.
  \item \textbf{Dirección de cara definitiva}: $\arctan2(-\Delta y, -\Delta x)$,
        signo consistente con el planificador.
  \item \textbf{Sin teletransporte}: \texttt{plan\_path\_only()} + camino 800\,ms
        antes de ejecutar.
  \item \textbf{Hydra corregido}: \texttt{scene\_dataset\_config\_file} sin prefijo
        \texttt{+} en ejecución y documentación de \texttt{menu.sh}.
  \item \textbf{Repositorio GitHub limpio}: \texttt{mobileclip\_blt.ts} (571\,MB)
        eliminado del historial con \texttt{git-filter-repo}; push operativo.
  \item \textbf{Bugs corregidos}: IndexError boundary (clamp en
        \texttt{plan\_path\_only()}) + umbral YOLOE subido a 0,40.
  \item \textbf{Fórmula de calidad}: implementada en \texttt{postprocess\_heatmap()};
        componentes ordenados por quality.
\end{enumerate}

\subsection*{14 de abril de 2026}

\begin{enumerate}
  \item \textbf{Reestructuración del documento de proyecto}: se redefine la
        finalidad del proyecto con foco en razonamiento por habitaciones.
  \item \textbf{HSSD confirmado como dataset principal}: MP3D queda como soporte
        secundario; LabelMe diferido.
  \item \textbf{Arquitectura room-aware diseñada}: 8 fases (A--H) con
        representación de estado por habitación, scoring de selección de room,
        acciones de alto nivel y rol estratégico del LLM.
  \item \textbf{Roadmap corregido}: la escalera M1--M5 se reemplaza por la
        arquitectura room-aware como prioridad central. Baselines y evaluación
        pasan a fase posterior.
  \item \textbf{Fase A implementada}: \texttt{SemanticSceneRoomProvider} construido
        con transformación completa Habitat$\rightarrow$grid (\texttt{hab\_to\_grid}).
        Polígonos rasterizados en espacio grid. Cada componente del heatmap anotado
        con su habitación. Navegación a habitación por nombre funcional (centroide
        directo, sin heatmap ni YOLOE). Verificación visual de llegada a habitación
        omitida en simulación (coordenadas exactas); se contempla reactivarla para
        robot real con localización imperfecta.
  \item \textbf{Parsing semántico por niveles}: identificado como trabajo futuro
        necesario. El parsing actual separa instrucciones compuestas en targets
        independientes; se requiere un parser que interprete habitación objetivo +
        objeto dentro de ella como una sola instrucción estructurada.
\end{enumerate}

\subsection*{15 de abril de 2026}

\begin{enumerate}
  \item \textbf{Fase B completada}: \texttt{RoomState} y \texttt{SearchState}
        implementados en \texttt{vlmaps/utils/search\_state.py}. Cada
        \texttt{RoomState} registra \texttt{explored\_ratio}, \texttt{objects\_seen},
        \texttt{candidates\_tried/confirmed}, \texttt{times\_visited},
        \texttt{target\_relevance} y \texttt{target\_found\_here}. Se construye una
        instancia por target al inicio de cada instrucción y se actualiza
        incrementalmente durante la navegación y la verificación YOLOE. Al final de
        cada instrucción se imprime un resumen completo del estado de búsqueda.

  \item \textbf{Fase C completada}: módulo \texttt{vlmaps/utils/room\_priors.py}.
        Implementa \texttt{compute\_room\_priors(query, known\_rooms,
        objects\_seen\_by\_room)} fusionando tres fuentes con pesos
        $a=0{,}5$, $b=0{,}3$, $c=0{,}2$: (a) LLM GPT-4o-mini con llamada
        estructurada al inicio de cada búsqueda de objeto; (b) tabla manual de
        $\sim$35 objetos con scores por tipo de habitación; (c) evidencia de
        co-ocurrencia basada en \texttt{objects\_seen} por habitación. El resultado
        normalizado se almacena en \texttt{RoomState.target\_relevance} y se
        muestra como traza de depuración antes de cada consulta.

  \item \textbf{Navegación solo a candidatos del heatmap}: eliminado el uso de
        \texttt{get\_standoff\_pos()} (que usaba \texttt{labeled\_map\_cropped}).
        Ahora el robot navega directamente al centroide del mejor componente de
        \texttt{kept\_components} (filtrado real por argmax + score $>$ 0{,}3).
        Esto evita que el robot viaje a zonas con señal VLMap pero sin activación
        real en el heatmap.

  \item \textbf{Filtrado de objetos navegables}: la lista de objetos mostrada antes
        de cada consulta usa el mismo filtro que \texttt{compute\_heatmap} (vóxeles
        con argmax ganador y score $> 0{,}3$, mínimo 10 vóxeles), garantizando
        coherencia entre lo que se muestra y lo que el sistema puede encontrar.

  \item \textbf{Correcciones de navegación}: (a) \texttt{n\_actions == 0} tratado
        como ``ya en el destino'' en vez de ``sin camino'', corrigiendo el caso
        post-llegada a habitación; (b) eliminado el bloque de reintento por paso
        estrecho (\textit{narrow-passage retry}) que generaba más falsos positivos
        que beneficios; (c) umbral YOLOE subido a $0{,}3$ y dilatación de obstáculos
        a 3 iteraciones ($\approx$15\,cm por lado).

  \item \textbf{Tipos canónicos de habitación (\textit{room\_type} vs
        \textit{room\_instance})}: HSSD codifica instancias repetidas de un mismo
        tipo de habitación con sufijos numéricos (\texttt{bathroom.001},
        \texttt{bathroom.002}, etc.). Se añade \texttt{canonical\_room\_type()} en
        \texttt{room\_priors.py} que elimina el sufijo \texttt{.NNN} mediante regex
        (\texttt{re.sub(r'\textbackslash.\textbackslash d+\$', '', s)}), y
        \texttt{compatible\_room\_types()} que devuelve el conjunto de tipos
        canónicos con prior $>0{,}05$. Estas funciones permiten razonar sobre
        familias de habitación independientemente del índice de instancia.

  \item \textbf{Comportamiento \textit{local-first} al navegar (\textit{room gate})}:
        nueva función \texttt{select\_best\_candidate(kept\_components, current\_room,
        query\_priors, tried\_centroids)} en \texttt{interactive\_object\_nav.py}. Si
        el robot ya se encuentra dentro de una instancia de habitación cuyo tipo
        canónico es compatible con el objeto buscado (e.g.\ robot en
        \texttt{bathroom.001}, objeto \texttt{sink}), el sistema prioriza los
        candidatos del heatmap situados en esa misma instancia, aplicando un bonus de
        $+0{,}35$ a su calidad efectiva. El cambio de habitación sólo se produce si
        (a) no quedan candidatos locales sin explorar, o (b) la calidad externa supera
        la local en más de un 40\% tras el bonus. El conjunto
        \texttt{\_tried\_centroids} en \texttt{SearchState} registra los centroides ya
        visitados para evitar ciclos.

  \item \textbf{Navegación con margen de seguridad aumentado}: la dilatación de
        obstáculos en \texttt{habitat\_lang\_robot.py} se eleva a 5 iteraciones
        ($\approx$25\,cm por lado). La función
        \texttt{select\_safe\_goal\_from\_path()} aumenta el radio mínimo de
        \textit{clearance} de 3 a 5 celdas y adopta estrategia de máxima holgura:
        reúne \emph{todos} los puntos del camino que cumplen el radio y selecciona el
        de mayor distancia al obstáculo más cercano (en vez del primero válido caminando
        hacia atrás). Si la pasada estricta no devuelve candidatos, una segunda pasada
        suaviza el umbral y vuelve a seleccionar por máxima holgura.

  \item \textbf{Métricas de seguridad por camino}: nueva función
        \texttt{\_compute\_path\_safety(path\_cells, obs\_map)} basada en
        \texttt{distance\_transform\_edt}. Registra longitud, holgura mínima y media,
        penalización acumulada por celdas con $\text{clearance} < 8$ y coste seguro
        $= \text{longitud} + 0{,}5 \cdot \text{penalización}$. Los valores se imprimen
        en trazas \texttt{[safety]} tras cada planificación, facilitando el diagnóstico
        de rutas peligrosas y la futura migración a ROS/HSR.

  \item \textbf{Corrección Bug 1 — comandos de habitación con validación real de
        llegada}: los comandos de tipo \texttt{kitchen}, \texttt{dining room},
        etc.\ fallaban porque el sistema navegaba al centroide geométrico del polígono
        de la habitación (a menudo sobre un obstáculo), el planificador lo desplazaba a
        un nodo visgraph cercano (frecuentemente en otra habitación) y después declaraba
        éxito sin comprobar dónde estaba realmente el robot. Se implementa
        \texttt{find\_reachable\_room\_goal()}, que enumera todas las celdas
        navegables del interior de la habitación objetivo, las ordena por holgura
        respecto a obstáculos y devuelve los top-$K$ candidatos seguros. Tras la
        ejecución del camino se comprueba la habitación real con
        \texttt{get\_room\_at\_cell()} y se valida mediante
        \texttt{\_room\_instance\_matches()}, que acepta tanto coincidencia exacta como
        de tipo canónico (\texttt{kitchen.001} satisface el comando \texttt{kitchen}).
        Si la primera meta falla, se reintenta con los candidatos alternativos antes de
        declarar fracaso. Solo se imprime ``Arrived at room X'' cuando la comprobación
        pasa.

  \item \textbf{Corrección Bug 2 — evidencia del heatmap domina en consultas
        directas}: para consultas de objeto con señal real en el heatmap (p.ej.\
        \texttt{sofa}), los priors de habitación dependían excesivamente de los
        priors lingüísticos (LLM + tabla manual), ignorando que los componentes del
        heatmap ya indicaban en qué habitación está el objeto. Además, nombres de
        habitación específicos de la escena como \texttt{tv} no coincidían con los
        fragmentos estándar (\texttt{living}) de la tabla manual. Se añade la tabla
        \texttt{\_SEMANTIC\_ALIASES} que mapea nombres atípicos a familias semánticas
        (\texttt{tv} $\to$ \texttt{living}, \texttt{laundryroom} $\to$
        \texttt{laundry}, etc.) y se aplica en \texttt{\_manual\_prior()} y
        \texttt{\_evidence\_prior()} mediante \texttt{\_normalize\_room\_for\_priors()}.
        Se añade \texttt{compute\_heatmap\_room\_evidence()} que agrega la calidad de
        componentes por instancia de habitación y normaliza a $[0,1]$. Se extiende
        \texttt{compute\_room\_priors()} con los parámetros opcionales
        \texttt{heatmap\_evidence} y \texttt{query\_type}: las consultas directas usan
        pesos $0{,}65 \cdot \text{heatmap} + 0{,}20 \cdot \text{manual} + 0{,}15
        \cdot \text{LLM}$, mientras que las indirectas mantienen
        $0{,}45 \cdot \text{LLM} + 0{,}30 \cdot \text{manual} + 0{,}25 \cdot
        \text{evidencia escena}$. En el bucle de navegación, tras anotar los
        componentes con su habitación, se recalculan los priors con evidencia real
        antes de seleccionar el candidato, de modo que \texttt{sofa} con tres
        componentes en \texttt{tv} produce un prior $>0{,}86$ para \texttt{tv}.

  \item \textbf{Corrección Bug 3 — puerta de seguridad bloquea ejecución de rutas
        inseguras}: las métricas de seguridad \texttt{[safety]} existían pero eran
        puramente informativas: el sistema ejecutaba igualmente rutas con
        \texttt{goal\_cl=0.0} o \texttt{min\_cl=0.0}. Se añaden umbrales duros
        configurables \texttt{\_MIN\_GOAL\_CLEARANCE=3{,}0} celdas y
        \texttt{\_MIN\_PATH\_CLEARANCE=1{,}0} celdas. Antes de ejecutar cualquier
        camino a un objeto, se evalúan dichos umbrales: si el goal o la ruta no los
        superan, el centroide se marca como intentado en \texttt{\_tried\_centroids}
        y se prueban hasta tres candidatos alternativos. Si ningún candidato
        supera los umbrales, la consulta se omite con mensaje explícito. La
        seguridad ya no es anotación, sino control de flujo.

  \item \textbf{Corrección Problema 1 — metas de aproximación seguras alrededor del
        centroide de componente}: el pipeline anterior planificaba hasta el centroide
        del heatmap (a menudo situado sobre un obstáculo), ejecutaba el camino completo
        y luego recorría el camino hacia atrás buscando un buen punto de observación. Si
        el centroide caía dentro de un obstáculo, el planificador visgraph producía rutas
        erróneas. Se añade \texttt{find\_safe\_candidate\_approach\_goals(centroid,
        safe\_obs\_map, robot\_pos)} que muestrea posiciones de aproximación en un anillo
        de $8$--$25$ celdas alrededor del centroide, filtra por navegabilidad en
        \texttt{\_safe\_obs\_map} y holgura mínima $\ge 3$ celdas, y ordena por holgura
        descendente y proximidad al robot. La meta de navegación se elige directamente
        del mejor candidato del anillo; el camino hacia atrás solo se activa como
        respaldo si el anillo queda vacío.

  \item \textbf{Corrección Problema 2 — planificador global atravesaba paredes}:
        causa raíz doble. (a) La dilatación de obstáculos con 5 iteraciones (bilateral
        $\approx10$ celdas) cerraba los pasillos estrechos de HSSD ($\sim$10 celdas de
        ancho) desconectando habitaciones en el visgraph. Corregido volviendo a
        3~iteraciones ($\approx6$ celdas). (b) El método de recuperación de goal en
        espacio de obstáculo usaba \texttt{G.closest\_point()} (holgura\,=\,0) con
        \texttt{get\_nearby\_position()} como fallback, que devuelve \texttt{None} si
        las cuatro diagonales están bloqueadas; \texttt{vg.Point(None[0])} genera
        coordenadas inválidas $(-1,\,96)$ y \texttt{G.shortest\_path} produce una línea
        recta a través de paredes. Ambas funciones se reemplazan por
        \texttt{\_snap\_to\_nearest\_free(point, obstacles, min\_clearance)} en
        \texttt{navigation\_utils.py}, que usa \texttt{distance\_transform\_edt} para
        encontrar la celda libre más cercana con holgura suficiente: el resultado está
        siempre garantizado dentro del espacio navegable del mapa del visgraph.
        Se añaden también registros de diagnóstico \texttt{[planner] Start navigable:
        \ldots{} Goal navigable: \ldots{}} y almacenamiento del mapa dilatado completo
        como \texttt{robot.\_safe\_obs\_map} para que el resto del sistema use la misma
        referencia que el visgraph. Las funciones
        \texttt{find\_reachable\_room\_goal()} y
        \texttt{find\_safe\_candidate\_approach\_goals()} se actualizan para usar
        \texttt{robot.\_safe\_obs\_map} en lugar del mapa de navmesh crudo, eliminando
        la discordancia entre validación de metas y planificador.

  \item \textbf{Corrección Problema 3 — escudo de colisiones local en ejecución}:
        \texttt{execute\_nav\_replay()} acepta ahora el parámetro opcional
        \texttt{dist\_map} (transformada de distancias del mapa seguro dilatado,
        precalculada una vez antes de cada ejecución). Antes de cada acción
        \texttt{move\_forward}, el escudo predice la siguiente celda de cuadrícula a
        partir de la posición actual \texttt{curr\_pos\_on\_map} y el ángulo
        \texttt{curr\_ang\_deg\_on\_map} del robot. Si la holgura en esa celda es
        inferior a $\texttt{STOP\_CL}\,{=}\,2{,}0$ celdas, la ejecución se interrumpe
        inmediatamente y se activa la recuperación; si la holgura está entre
        $\texttt{STOP\_CL}$ y $\texttt{SLOW\_CL}\,{=}\,5{,}0$ celdas se emite un aviso
        sin detener el movimiento. El escudo opera sobre el mismo mapa dilatado que el
        visgraph, de modo que «celda segura» tiene idéntico significado en planificación
        y en ejecución.

  \item \textbf{Navegación one-shot — eliminados reintentos y recuperación}:
        el escudo de colisiones usaba el mapa dilatado (\texttt{\_safe\_obs\_map})
        para calcular la transformada de distancias, lo que hacía que el centro de
        un pasillo de $\sim10$ celdas tuviera clearance $\approx 2$ en el mapa
        dilatado y disparara paradas falsas al cruzar puertas. El escudo pasa a usar
        \texttt{robot.map.obstacles\_map} (mapa crudo sin dilatar), que refleja el
        espacio físico real; los umbrales se ajustan a
        $\texttt{STOP\_CL}=1{,}0$ y $\texttt{SLOW\_CL}=2{,}5$ celdas.
        Adicionalmente, se elimina la llamada a \texttt{nav\_recovery\_and\_replan()}
        tras un fallo de ejecución y el bucle de reintentos con metas alternativas
        en comandos de habitación. La navegación es ahora estrictamente one-shot:
        se planifica una vez, se ejecuta una vez y se reporta el estado real de la
        habitación; si la ejecución se interrumpe, el robot continúa desde la posición
        actual sin ruta nueva.
\end{enumerate}

% =============================================================================
\subsection*{16 de abril de 2026}
% =============================================================================

\begin{enumerate}

  \item \textbf{Separación explícita entre inflación de planificación y escudo
        de ejecución}: el planificador global utiliza el mapa dilatado
        (\texttt{\_safe\_obs\_map}, 3~iteraciones de dilatación binaria,
        $\approx$15\,cm por lado) para trazar rutas con suficiente margen.
        El escudo de ejecución opera exclusivamente sobre la transformada de
        distancias del mapa crudo (\texttt{robot.map.obstacles\_map}), sin
        dilatar. Esto garantiza que los pasillos estrechos de HSSD
        ($\sim$10~celdas de ancho, clearance real~$\approx$5~celdas) no sean
        bloqueados falsamente por el escudo. La constante
        \texttt{\_PLAN\_DILATION\_ITERS = 3} se declara de forma explícita en
        \texttt{habitat\_lang\_robot.py} y se imprime al inicio de la sesión
        junto con el mensaje: «\textit{global planner only; execution shield
        uses raw map}».

  \item \textbf{Modo de travesía de puertas (\textit{doorway mode})}: antes de
        cada \texttt{move\_forward}, el escudo mide las clearances
        perpendiculares a la derecha y a la izquierda del robot en la pose
        actual (\texttt{\_side\_l}, \texttt{\_side\_r}). Si ambas caen por
        debajo de \texttt{DOORWAY\_SIDE\_TH}~$= 5{,}0$~celdas y la clearance
        frontal supera \texttt{DOORWAY\_FRONT\_MIN}~$= 0{,}5$~celdas, se
        activa el modo de paso estrecho. En ese modo el umbral de parada dura
        se relaja a \texttt{STOP\_CL\_DOOR}~$= 0{,}5$~celdas (frente a
        1{,}0 en espacio abierto) y cada paso se registra como
        «\textit{cautious doorway traversal}». La transición
        ON/OFF se imprime explícitamente. El planificador puede enrutar
        libremente por puertas válidas; el escudo deja cruzarlas sin falsas
        paradas.

  \item \textbf{Escudo sensible a la huella completa del robot}: el chequeo
        anterior solo evaluaba un punto (celda frontal prevista). Ahora se
        evalúan tres puntos de la huella en la pose predicha:
        \textit{front-center}, \textit{front-left} y \textit{front-right}
        (desplazados lateralmente \texttt{ROBOT\_HW}~$= 2$~celdas respecto
        al eje de avance). Se registra por separado
        \texttt{front clearance}, \texttt{left clearance},
        \texttt{right clearance} y \texttt{footprint min clearance}. En
        espacio abierto, si algún lateral cae por debajo de
        \texttt{SIDE\_STOP\_CL}~$= 1{,}5$~celdas, el paso se bloquea. En
        modo puerta, el criterio de bloqueo es solo el mínimo de la huella
        frente a \texttt{DOORWAY\_FRONT\_MIN}.

  \item \textbf{Puerta de entrada a habitación para confirmación YOLOE (\textit{room-entry gate})}: hasta ahora el sistema podía declarar éxito
        inmediatamente si YOLOE detectaba el objeto desde la puerta antes de
        haber cruzado el umbral. Se añade una comprobación en los stages 0
        y~1 de YOLOE: si el robot está en una habitación distinta de la que
        contiene el componente del heatmap objetivo
        (\texttt{best\_comp["room"]}) y la distancia al centroide del
        candidato supera 20~celdas ($\approx$1\,m), la detección se marca
        como \textit{tentativa} (\texttt{\_yoloe\_confirmed = False}) y la
        búsqueda continúa con el stage siguiente. Solo se acepta como éxito
        definitivo si el robot ha entrado en la habitación correcta o se
        encuentra a menos de 20~celdas del candidato. Se registra en consola
        «\textit{Object visible but room not yet entered → tentative only}» o
        «\textit{Room entry confirmed}» según el caso.

  \item \textbf{Centrado de trayectorias en pasos estrechos mediante eje medial local}:
        el visgraph seguía generando caminos geométricamente válidos pero
        sesgados hacia uno de los lados del estrechamiento, lo que hacía que el
        escudo de ejecución bloquease pasos por clearance lateral insuficiente
        aun estando en modo de paso estrecho. El primer intento de corrección
        solo recentraba vértices ya existentes y no bastaba cuando el
        estrechamiento caía en mitad de un segmento largo del visgraph. Se
        añade un postprocesado
        \texttt{center\_path\_on\_medial\_axis()} en
        \texttt{vlmaps/utils/navigation\_utils.py}. Tras obtener la ruta del
        visgraph, se rasteriza cada segmento y, si detecta una franja con
        clearance $<4{,}0$ celdas, inserta un waypoint auxiliar en la zona más
        estrecha. Ese waypoint, junto con los intermedios originales, se
        desplaza hacia la celda de mayor holgura dentro de un disco local de
        radio 6, usando \texttt{distance\_transform\_edt} del mapa
        \emph{sin} dilatar. El candidato solo se acepta si mantiene línea de
        visión libre con los waypoints vecino anterior y siguiente sobre el
        mapa de planificación dilatado, evitando atajos que corten esquinas o
        crucen paredes. Después se aplica un suavizado conservador de 3 puntos
        y se conserva la meta final sin desplazar. El resultado es que la ruta
        cruza por el centro del paso incluso cuando el cuello de botella no
        coincide con un vértice del visgraph, y el escudo vuelve a comportarse
        como mecanismo de seguridad, no como corrector sistemático de
        trayectorias.

  \item \textbf{Modo de paso estrecho activado por geometría local, no solo por
        simetría lateral}: en ejecución real aparecía un caso límite al entrar
        oblicuamente en un marco o cuello de botella: un lateral del footprint
        caía por debajo de 1{,}5 celdas mientras el otro seguía aún en espacio
        relativamente abierto, de modo que el detector antiguo no activaba el
        modo de paso estrecho y el escudo bloqueaba el avance con la regla de
        espacio abierto. Se refuerza la detección en
        \texttt{execute\_nav\_replay()} con tres señales:
        (1) estrechamiento en la pose actual, (2) estrechamiento en la huella
        de la pose siguiente y (3) entrada asimétrica, cuando un lateral ya es
        crítico pero la clearance frontal sigue siendo cruzable. Además se
        añade histéresis de salida para mantener el modo activo mientras el
        robot permanezca dentro del cuello de botella. El modo pasa a modelar
        \emph{pasos estrechos cruzables} en general, no solo puertas
        perfectamente simétricas entre habitaciones.

  \item \textbf{Sustitución del \textit{waypoint chasing} por un seguidor de
        polilínea con \textit{lookahead}}: el aumento de waypoints auxiliares
        para centrar la trayectoria en pasos estrechos mejoraba la geometría,
        pero degradaba la ejecución con cientos de micro-correcciones de rumbo,
        porque el controlador discreto original perseguía cada subgoal por
        separado con cuantización fija de 5° y 0{,}1\,m. Se densifica ahora la
        ruta planificada a nivel de celda mediante rasterización de segmentos y
        se ejecuta con un \textit{path follower} sobre esa polilínea densa. En
        cada iteración, el robot proyecta su pose sobre la ruta, selecciona un
        objetivo adelantado (\textit{lookahead}) varios pasos por delante y
        solo aplica la primera acción discreta necesaria para aproximarse a ese
        punto. El \textit{lookahead} se reduce automáticamente dentro de pasos
        estrechos y se mantiene más largo en espacio abierto. Además, las
        métricas de seguridad y la selección de metas de observación pasan a
        evaluarse sobre la misma polilínea densificada, de modo que
        planificación, validación y ejecución comparten por fin la misma
        representación geométrica del camino.

  \item \textbf{\textit{Lookahead} recortado por visibilidad local del camino}:
        el primer seguidor con \textit{lookahead} seguía teniendo un fallo
        importante en pasillos con giro o marcos oblicuos: podía elegir como
        objetivo un punto varios pasos por delante que ya estaba al otro lado
        de una esquina no visible desde la pose actual. Como el controlador
        discreto solo orienta hacia ese objetivo y aplica la primera acción,
        el robot intentaba cortar la esquina y acababa proponiendo un avance
        con \texttt{front clearance = 0.0}. Se corrige haciendo que el
        \textit{path follower} no use directamente el \textit{lookahead}
        preferido, sino el punto más lejano de la polilínea densa que siga
        siendo alcanzable con línea de visión libre sobre el mapa de
        planificación dilatado. Así, dentro de cuellos de botella y giros
        cerrados, el seguidor deja de apuntar ``a través de la pared'' y se
        mantiene sobre el tramo realmente visible del camino.

  \item \textbf{Etiquetas de llegada coherentes con el resultado real de la
        navegación}: en comandos de habitación se estaba mostrando
        ``\textit{Arrived: ...}'' inmediatamente después de que terminara la
        ejecución del follower, incluso si el robot se había detenido antes del
        umbral y la comprobación final de habitación fallaba. Esto inducía a
        error durante la depuración visual. Se corrige la interfaz para que el
        overlay solo muestre ``Arrived'' cuando el robot ha alcanzado de verdad
        la habitación objetivo; en caso contrario, ahora se muestra
        ``Not reached'' en navegación a habitaciones o ``Stopped before goal''
        en aproximaciones a objetos.

  \item \textbf{La referencia de control vuelve a ser la ruta densa, no solo
        la polilínea compactada}: la versión que seguía únicamente la
        polilínea geométrica simplificada seguía permitiendo atajos laterales
        del controlador discreto. En la práctica, el robot podía ``llegar más
        o menos'' al final del segmento pero desviándose del eje real del
        camino, saliéndose hacia paredes o quedándose en la habitación vecina
        cuando el objetivo estaba cerca de una frontera semántica. Se
        reintroduce la trayectoria densa rasterizada como referencia de
        control, pero sin perseguir cada celda individualmente: el follower
        mantiene \textit{lookahead} sobre esa ruta densa y la usa para medir
        progreso y error lateral. La polilínea compactada queda como
        representación del plan, mientras que la ejecución se ancla a la traza
        detallada que realmente debe seguir el robot.

  \item \textbf{Banda muerta angular para eliminar micro-correcciones de 5°}:
        incluso después de volver a seguir la polilínea geométrica, el
        controlador discreto seguía tendiendo a consumir pequeños giros antes
        de casi cada avance cuando el error angular hacia el siguiente objetivo
        era muy bajo. Eso no cambiaba de forma significativa la seguridad, pero
        sí alargaba innecesariamente la ejecución en pasos estrechos. Se añade
        una tolerancia angular en el follower: si el objetivo inmediato ya está
        alineado dentro de un margen moderado, se prioriza
        \texttt{move\_forward} en lugar de corregir rumbo con otro giro
        discreto. La tolerancia es más amplia en zonas estrechas que en espacio
        abierto. Con ello, se reduce el número de pasos y el patrón de
        micro-corrección sin introducir nuevos modos especiales de paso por
        puertas.

  \item \textbf{El shield preventivo se elimina de la ejecución paso a paso}:
        el análisis fino mostró que el escudo local estaba mezclando dos
        nociones incompatibles: clearance sobre un mapa ya dilatado por el
        radio del agente y probes adicionales alrededor del cuerpo. Eso
        generaba falsos bloqueos en puertas y convertía la seguridad en el
        mecanismo principal de decisión, cuando debería ser solo una red de
        protección. La arquitectura actual elimina el \textit{pre-step shield}
        y delega la colisión física en la navmesh de Habitat. La robustez pasa
        a depender de dos piezas más limpias: planificación \textit{clearance-
        aware} y seguimiento fiel del camino.

  \item \textbf{Watchdog de no-progreso real en el follower}: además del fallo
        semántico del shield, aparecía un segundo problema en ejecución: el
        robot podía quedarse vivo cientos de iteraciones repitiendo
        prácticamente la misma pose y agotando el \textit{step budget}. Se
        añade un detector de no-progreso que vigila tanto la celda actual como
        el índice efectivo de progreso sobre la polilínea. Si durante varias
        iteraciones no cambia ni la celda útil ni el progreso útil, la
        ejecución se aborta pronto. También se detecta el caso en que el
        controlador no produce acciones válidas repetidamente para el mismo
        target local. La versión inicial de este watchdog abortaba
        incorrectamente durante la fase de giro en sitio, porque tomaba
        ``misma celda'' como sinónimo de ``sin progreso'' incluso cuando la
        orientación seguía cambiando de forma útil hacia el siguiente tramo del
        camino. Se corrige para que una rotación real también cuente como
        progreso temporal del follower. Con ello, los fallos pasan a ser
        rápidos y diagnosticables en lugar de consumir cientos de pasos
        inmóviles o abortar antes de empezar a avanzar.

  \item \textbf{Reenganche explícito al camino cuando el robot se sale del
        eje}: incluso sin \textit{shield}, el follower seguía pudiendo
        comportarse de forma errática si la banda muerta angular le dejaba
        avanzar cuando ya se había separado lateralmente del trazado. Se añade
        una lógica de \textit{off-track reattachment}: la distancia del robot a
        la ruta densa se mide en cada iteración y, si supera un umbral, el
        \textit{lookahead} se acorta automáticamente para volver a engancharse
        al eje del camino antes de permitir más avance agresivo. Además, la
        banda muerta angular solo puede forzar \texttt{move\_forward} cuando el
        error lateral sigue siendo pequeño.

  \item \textbf{Llegada más estricta en navegación a habitaciones}: la lógica
        anterior consideraba alcanzado el objetivo si el robot quedaba a pocas
        celdas del final del camino. Eso era suficiente para aproximaciones a
        objetos, pero no para comandos de habitación cerca de fronteras
        semánticas: el robot podía detenerse a 1--2 celdas del goal y seguir
        clasificado en la estancia vecina. Se introduce una tolerancia de
        llegada más estricta para navegación a habitaciones, de modo que el
        follower no dé por completado el camino hasta estar realmente encima
        del objetivo final o prácticamente superpuesto a él.

  \item \textbf{Deadband angular menos restrictiva en el follower}: la primera
        versión de la banda muerta angular seguía siendo demasiado conservadora
        porque solo permitía sustituir un giro por un avance si el segundo
        elemento del plan local del controlador ya era \texttt{move\_forward}.
        En la práctica, eso dejaba fuera muchos casos útiles en los que el
        error angular ya era pequeño pero el controlador aún proponía varios
        giros discretos consecutivos. Se elimina esa condición adicional: ahora
        el follower prioriza \texttt{move\_forward} siempre que el error angular
        esté dentro de la tolerancia y la distancia al objetivo inmediato siga
        justificando el avance. Esto reduce el patrón
        ``giro-giro-giro-avance'' en puertas y pasillos angostos sin cambiar la
        arquitectura general del controlador.

\end{enumerate}

% =============================================================================
\subsection*{20 de abril de 2026}
% =============================================================================

\begin{enumerate}

  \item \textbf{Validación en vivo dentro del contenedor de simulación}: en
        lugar de seguir depurando solo con logs pegados manualmente, se pasa a
        ejecutar pruebas reales dentro de \texttt{tfg-sim} sobre la escena
        HSSD usada durante el desarrollo (\texttt{scene\_id=1}). Esto permitió
        comprobar de forma reproducible tres consultas representativas:
        navegación a habitación (\texttt{kitchen}), búsqueda de objeto con
        ambigüedad semántica local (\texttt{sofa}) y búsqueda de objeto con
        varios candidatos plausibles (\texttt{sink}). La prueba en vivo mostró
        que \texttt{kitchen} seguía fallando aunque el \textit{shield} ya no
        bloqueaba puertas, mientras que \texttt{sofa} y \texttt{sink} sí
        llegaban a confirmarse con YOLOE. Con ello quedó claro que el problema
        residual principal ya no estaba en la planificación ni en la
        colisión preventiva, sino en el seguimiento local del camino.

  \item \textbf{Corrección del bucle \texttt{preview produced no action}}:
        el fallo real de \texttt{kitchen} no era una colisión contra paredes,
        sino un colapso del follower cerca de un tramo concreto del camino,
        donde el objetivo local quedaba tan cerca que el controlador discreto
        devolvía repetidamente una previsualización vacía
        (\texttt{progress=58, target=59}). El follower seguía reenganchándose a
        ese mismo índice una y otra vez y acababa abortando sin salir del
        cuello. Se corrige en dos pasos. Primero, el objetivo visible de
        \textit{lookahead} ya no puede ser arbitrariamente cercano salvo cuando
        el robot está realmente desviado o casi en la meta. Segundo, cuando un
        target local sigue produciendo una previsualización vacía, el sistema
        eleva explícitamente un \textit{progress floor} y fuerza al follower a
        no volver a anclarse en ese mismo punto de la ruta. Con este cambio, la
        consulta \texttt{kitchen} deja de abortar en el mismo progreso y
        completa la trayectoria hasta quedar clasificada correctamente dentro de
        la cocina.

  \item \textbf{La calidad semántica del heatmap deja de premiar blobs
        minúsculos}: el caso de \texttt{sofa} reveló otro problema de calidad.
        La fórmula anterior para puntuar componentes del heatmap aún permitía
        que un blob muy pequeño pero extremadamente denso compitiera con una
        región grande correspondiente al sofá real. En la práctica, el robot
        podía priorizar un sillón o una activación residual cercana solo porque
        el componente era compacto, aunque el soporte espacial del sofá fuese
        mucho mayor. Se reformula la \textit{quality} de cada componente para
        favorecer de forma más explícita el soporte semántico real
        (\texttt{sum\_val}) y el tamaño útil del componente, reduciendo el peso
        relativo de activaciones pequeñas. Tras este cambio, en la consulta
        \texttt{sofa} el componente grande asociado al sofá pasa a quedar por
        delante de los blobs minúsculos, evitando que el robot ``vaya al
        sillón'' salvo que semánticamente sea de verdad el mejor candidato.

  \item \textbf{Overlay y visualización más fluidos sin cambiar de backend}:
        también se revisa la parte puramente visual del sistema. El overlay de
        navegación mostraba etiquetas del tipo \texttt{[64/366] -> sofa}, pero
        ese total era poco informativo para el usuario y además enfatizaba lo
        largo que era el seguimiento discreto. Se simplifica el texto en
        pantalla para mostrar solo el paso actual (\texttt{[64] -> sofa}). A
        nivel de fluidez, no se migra al visualizador 3D de Habitat, porque
        eso aumentaría la complejidad sin resolver el cuello principal. En su
        lugar, se mantiene OpenCV y se hace más eficiente: la vista en primera
        persona se sigue actualizando en cada acción, pero el mapa cenital se
        redibuja con una cadencia menor y el retardo explícito por paso se
        reduce prácticamente a cero. Así, la reproducción visual resulta mucho
        más fluida sin penalizar perceptiblemente el rendimiento de la
        simulación ni complicar la arquitectura de depuración actual.

  \item \textbf{Fase D implementada en el pipeline real}: el razonamiento por
        habitaciones deja de ser una traza auxiliar y pasa a intervenir en la
        decisión. Se añade una función \texttt{select\_best\_room()} que, a
        partir de \texttt{SearchState}, la evidencia del heatmap y la posición
        actual del robot, calcula un \texttt{room\_score} por habitación con
        cuatro señales ya disponibles en el sistema: prior semántico, evidencia
        agregada del heatmap, coste BFS sobre el mapa libre y penalización por
        intentos fallidos previos en esa habitación. La habitación ganadora se
        usa inmediatamente para filtrar \texttt{kept\_components} antes del
        ranking fino de candidatos, de modo que el flujo deja de ser
        puramente global. La implementación actual utiliza una proxy práctica
        de \textit{unexplored} basada en visitas y fallos acumulados, y deja
        el seguimiento de cobertura detallada para la futura Fase~F.

  \item \textbf{Ajustes del flujo de búsqueda visual y preparación de tests por
        lote}: tras estabilizar la Fase~D se corrigen varios detalles del flujo
        objeto$\rightarrow$detección. Primero, la selección de goals de
        aproximación deja de aceptar puntos que, aun teniendo clearance y
        perteneciendo a la misma habitación, quedarían geométricamente al otro
        lado de una pared respecto al centroid del candidato; ahora se exige
        además línea libre sobre el mapa de planificación. Segundo, las
        consultas directas a muebles presentes en la escena (por ejemplo
        \texttt{sofa} o \texttt{sink}) dejan de pasar por el \textit{room
        selector}: en ese modo se fuerza una política \textit{nearest-first}
        sobre los candidatos del heatmap y el razonamiento por habitaciones se
        reserva para consultas indirectas o ambiguas. Tercero, cuando la query
        sí activa selección de habitación, se añade un \textit{room stage}
        explícito: el robot navega primero a un goal interior seguro de la sala
        elegida y solo después realiza el acercamiento al candidato concreto,
        evitando verificar desde el centro de una habitación equivocada.
        Cuarto, la confirmación YOLOE pasa a congelar tanto la vista RGB
        anotada como el punto objetivo sobre el mapa cenital, dejando una marca
        persistente de la detección positiva. Finalmente, se preparan scripts
        de \textit{batch testing} para generar automáticamente una batería de
        consultas con todas las habitaciones suficientemente navegables y con
        todos los objetos presentes en la escena, dejando el runtime real de
        esas pruebas para la siguiente ronda de validación manual.

\end{enumerate}

% =============================================================================
\subsection*{21 de abril de 2026}
% =============================================================================

\begin{enumerate}

  \item \textbf{Fase F consolidada y accesible para validación manual}: se da
        por cerrada la capa de acciones de alto nivel introducida el día
        anterior y se revisa explícitamente su alcance real. La implementación
        no modifica todavía el comportamiento \emph{runtime} del robot dentro
        de \texttt{interactive\_object\_nav.py}; su función es preparar la
        transición desde una política \emph{inline} a una política explícita y
        tipada. Quedan definidos cinco tipos de acción
        (\texttt{go\_to\_room}, \texttt{explore\_room},
        \texttt{inspect\_candidate}, \texttt{verify\_target} y
        \texttt{done}), su validación estructural, su serialización JSON para
        la futura capa LLM y un selector de \emph{frontier} intra-habitación
        para \texttt{explore\_room}. En paralelo, \texttt{SearchState} amplía
        su papel como memoria del episodio con un registro ordenado de acciones
        y un rastro de celdas visitadas, que servirán como base del ejecutor de
        la siguiente fase.

  \item \textbf{Pruebas unitarias y acceso desde el menú del contenedor}: para
        que la validación de Fase~F no dependa de Habitat, de la GPU ni de una
        escena concreta, se mantiene una suite \texttt{pytest} autocontenida
        (\texttt{tests/test\_phase\_f\_actions.py}) y, además, se expone su
        ejecución desde el menú interactivo del contenedor mediante una nueva
        opción \texttt{p}. Esta entrada lanza la batería de pruebas de la capa
        de política y muestra de forma inmediata si la representación de
        acciones, el \textit{round-trip} JSON, el registro extendido de
        \texttt{SearchState} y el \textit{frontier picker} siguen siendo
        consistentes. El objetivo aquí no es medir navegación real, sino
        garantizar que la interfaz de alto nivel sobre la que se montará la
        siguiente fase permanece estable antes de conectarla al pipeline
        operativo.

  \item \textbf{Corrección del lanzador de \textit{batch testing}}: al ejecutar
        la opción \texttt{t} del menú con dataset HSSD, el script
        \texttt{tools/nav\_batch\_queries.py} fallaba con
        \texttt{TypeError: string indices must be integers} en
        \texttt{HabitatLanguageRobot.\_\_init\_\_}. La causa era que el script
        cargaba la configuración con \texttt{OmegaConf.load}, que no resuelve
        la lista \texttt{defaults} de Hydra: el campo \texttt{data\_paths}
        quedaba como la cadena literal \texttt{"default"} en lugar del
        diccionario con \texttt{habitat\_scene\_dir}/\texttt{vlmaps\_data\_dir}.
        Se sustituye la carga manual por \texttt{hydra.compose} con
        \texttt{initialize\_config\_dir}, pasando los \emph{overrides}
        (\texttt{scene\_id}, \texttt{dataset\_type},
        \texttt{scene\_dataset\_config\_file}, \texttt{data\_paths}) por la
        vía oficial de Hydra. Tras el arreglo, el script genera
        \texttt{queries.txt} correctamente y el batch llega a lanzar
        \texttt{interactive\_object\_nav.py}. Queda registrado además que un
        \texttt{scene\_id} fuera de rango (p.\,ej.\ teclear \texttt{11} cuando
        solo existen \texttt{0} y \texttt{1}) provoca un \texttt{IndexError}
        en \texttt{vlmaps\_data\_save\_dirs[scene\_id]}; es una limitación
        de UX del menú, no un bug del pipeline, y se considerará añadir una
        validación previa.

  \item \textbf{Primer corte funcional del executor alternativo}: se materializa
        por fin el apartado ``baseline vs.\ executor'' con una implementación
        inicial real. Se mantiene \texttt{interactive\_object\_nav.py} como
        \emph{baseline} intacto y, en paralelo, se añade
        \texttt{application/interactive\_object\_nav\_executor.py} como nuevo
        entrypoint junto con \texttt{vlmaps/policy/executor.py}. El nuevo
        módulo define un \texttt{ExecutorContext} y un despachador
        \texttt{execute\_action(...)} que traduce las acciones tipadas de
        Fase~F (\texttt{go\_to\_room}, \texttt{explore\_room},
        \texttt{inspect\_candidate}, \texttt{verify\_target},
        \texttt{done}) a comportamiento real reutilizando las primitivas ya
        validadas del baseline: \textit{room staging}, generación de
        \emph{approach goals}, \textit{path replay follower}, chequeos YOLOE
        en llegada y tras giro, \textit{scan}~$360^\circ$ y centrado visual.
        La política de alto nivel se mantiene coherente con el estado del
        proyecto: si la query es una habitación, se genera
        \texttt{go\_to\_room} $\rightarrow$ \texttt{done}; si es un mueble
        directo presente en la escena, se mantiene \textit{nearest-first}
        sin selector de habitación; y para consultas ambiguas o indirectas se
        conserva el flujo jerárquico ``habitación primero, candidato después''.
        Además, el menú del contenedor incorpora una nueva opción
        \texttt{e) Interactive executor navigation}, que permite lanzar esta
        variante sin tocar comandos manuales. La parte \emph{pendiente} ya no
        es construir el ejecutor, sino compararlo sistemáticamente frente al
        baseline con las métricas de Fase~H y con volcados lado a lado del
        heatmap baseline y del heatmap procesado \emph{room-aware}.

  \item \textbf{Resolución explícita de instancias de habitación en el parser y
        en la comprobación de llegada}: se corrige un fallo funcional relevante
        del flujo de navegación por habitaciones. Hasta este punto, órdenes como
        \texttt{bathroom.001} o expresiones más naturales como
        \texttt{bathroom 1} terminaban colapsándose sobre la habitación base
        \texttt{bathroom}, porque el pipeline solo distinguía familias
        semánticas y no instancias concretas. Se introduce un resolvedor
        explícito de nombres de habitación en \texttt{room\_provider.py}, capaz
        de mapear alias amigables
        (\texttt{bathroom 1}, \texttt{bathroom one}, \texttt{bedroom 2}) a sus
        etiquetas exactas de escena (\texttt{bathroom.001},
        \texttt{bedroom.002}, etc.), manteniendo a la vez el comportamiento
        esperado para comandos base como \texttt{bathroom} o \texttt{bedroom},
        que siguen prefiriendo la habitación principal cuando existe. Esta
        resolución se conecta tanto al \emph{baseline} como al nuevo
        \emph{executor}: tras el parse LLM se preservan menciones explícitas de
        instancia presentes en la instrucción original, la selección de
        \emph{room goals} pasa a usar la instancia resuelta y la condición de
        llegada deja de aceptar que \texttt{bathroom} satisfaga
        \texttt{bathroom.001}. Se añade además una batería unitaria nueva
        (\texttt{tests/test\_room\_instance\_resolution.py}) que valida la
        distinción entre habitación base e instancia numerada, junto con la
        compatibilidad con el resto de la capa de acciones de Fase~F.

  \item \textbf{Batería de comparación baseline vs.\ executor disponible desde
        el menú}: el apartado extra deja de ser solo una propuesta en papel y
        pasa a tener una infraestructura operativa para lanzar comparaciones
        controladas. Se añade al menú del contenedor la opción
        \texttt{v) Compare --- baseline vs executor}, que genera un lote común
        de \emph{queries}, ejecuta ambas variantes del navegador sobre ese mismo
        fichero de entrada y resume el resultado en \texttt{CSV} y
        \texttt{markdown}. La implementación queda repartida entre
        \texttt{tools/nav\_batch\_queries.py} (generación del lote),
        \texttt{tools/run\_hssd\_nav\_compare.sh} (orquestación host-side) y
        \texttt{tools/compare\_nav\_runs.py} (resumen query-a-query de
        detecciones, habitación final, abortos y número de actualizaciones de
        pose). No sustituye todavía a la evaluación de Fase~H, pero elimina el
        trabajo manual repetitivo de abrir y comparar logs a mano.

  \item \textbf{Primer corte funcional de Fase~G sobre el executor}: se añade
        \texttt{vlmaps/policy/strategic\_policy.py} como capa de decisión
        estratégica sobre estado estructurado. Esta pieza construye un
        \emph{snapshot} del episodio (prior de habitaciones, evidencia del
        heatmap, shortlist de candidatos, frontiers sugeridos y últimas
        acciones), consulta opcionalmente a un LLM para pedir la siguiente
        \texttt{Action} en formato JSON y valida esa propuesta antes de
        ejecutarla. Si el LLM no está disponible o devuelve una acción
        inconsistente, el sistema cae automáticamente a una política
        heurística determinista, lo que evita romper el flujo ya validado del
        executor. A nivel de integración, \texttt{interactive\_object\_nav\_executor.py}
        deja de preconstruir un plan fijo y pasa a operar como bucle
        ``elegir siguiente acción $\rightarrow$ ejecutar acción'', con modo de
        política seleccionable desde el menú (\texttt{heuristic},
        \texttt{hybrid}, \texttt{llm}). La validación unitaria queda cubierta
        por \texttt{tests/test\_phase\_g\_strategic\_policy.py}.

\end{enumerate}

% =============================================================================
\subsection*{22 de abril de 2026}
% =============================================================================

\begin{enumerate}

  \item \textbf{El \textit{testing} deja de ser una colección de utilidades y
        pasa a ser un flujo único}: se reorganiza por completo el submenú
        \texttt{t} del contenedor para que muestre solo tres opciones visibles:
        \texttt{Generate test set}, \texttt{Compare full 2x2 pipeline} y
        \texttt{Heatmap-only offline analysis}. Las antiguas utilidades
        intermedias (batch baseline aislado, comparación binaria baseline vs.\
        executor, runner JSONL suelto y tests unitarios de F/G) siguen
        existiendo como piezas internas o de desarrollo, pero dejan de ser el
        flujo principal del usuario. La intención es que la evaluación pase a
        ser reproducible y autoexplicativa: elegir escena(s), pulsar, esperar y
        recoger resultados en una carpeta estándar.

  \item \textbf{Selección explícita del tipo de heatmap en runtime}: hasta esta
        fecha, baseline y executor compartían implícitamente
        \texttt{compute\_heatmap()}, lo que impedía montar una comparación 2x2
        limpia sin duplicar entrypoints. Se introduce
        \texttt{VLMAPS\_HEATMAP\_MODE} con dos modos explícitos:
        \texttt{baseline} (heatmap crudo tras proyección 2D + decay, sin
        postproceso) y \texttt{postprocessed} (pipeline actual con
        \texttt{postprocess\_heatmap}). El baseline y el executor imprimen el
        modo activo al arrancar y el valor queda recogido en el
        \texttt{manifest.json} de cada corrida. Además, ambos entrypoints
        emiten ahora al final de cada query un resumen estructurado
        \texttt{[eval-summary] \{...\}} con habitación final, historial de
        habitaciones, número de candidatos inspeccionados/confirmados, rastro
        de acciones y recuento de celdas visitadas. Esta capa permite
        calcular métricas de Fase~H sin depender de \emph{regex} frágiles sobre
        logs textuales.

  \item \textbf{Runner online 2x2 y agregador de métricas completos}: se añaden
        dos herramientas nuevas. La primera,
        \texttt{tools/run\_full\_2x2\_eval.py}, coordina de forma automática las
        cuatro combinaciones \texttt{Ob\_Hb}, \texttt{Oe\_Hb},
        \texttt{Ob\_Hp} y \texttt{Oe\_Hp}, reutilizando
        \texttt{tools/run\_nav\_eval.py} como runner por método. La segunda,
        \texttt{tools/aggregate\_full\_2x2\_eval.py}, consume los cuatro
        manifiestos y genera dos vistas listas para inspección:
        \texttt{compare\_full.csv}/\texttt{.md} (una fila por query) y
        \texttt{aggregate\_metrics.csv}/\texttt{.md} (una fila por método y
        \emph{slice}). Las salidas se agrupan bajo una estructura homogénea:
        \texttt{/workspace/results/eval\_runs/\{timestamp\}/queries/},
        \texttt{pipeline\_2x2/\{scene\}/Ob\_Hb}, \ldots,
        junto con un \texttt{config.json} raíz que describe escenas, semilla,
        política del executor, rutas de queries y métodos corridos.

  \item \textbf{Métricas online ya instrumentadas y \emph{slicing} por
        \texttt{query\_type}/\texttt{tag}}: el agregador 2x2 calcula ya de forma
        reproducible SR, Object~SR, Room~SR, CFR, CT2R,
        Rooms~Before~Success, Inspections~Before~Success, Wrong~Visits,
        Mean~Pose~Updates, Early~Stop~Rate, Final~Room~Accuracy,
        Preview~No~Action~Rate, Path~Stopped~Early~Rate,
        Mean~Executor~Actions y Mean~Room~Transitions. No todas las métricas
        tienen aún el mismo grado de robustez: por ejemplo,
        \texttt{Inspections Before Success} se aproxima como
        \texttt{candidates\_tried - candidates\_confirmed} cuando la query
        termina en éxito, y \texttt{Rooms Before Success} se interpreta como el
        índice de la primera habitación esperada dentro de \texttt{visit\_history}.
        Estas aproximaciones quedan explícitas en el \texttt{Markdown} generado,
        en vez de presentarlas como magnitudes más precisas de lo que el log
        realmente permite.

  \item \textbf{La evaluación offline de heatmaps gana agregación por slices y
        envoltorio estándar}: \texttt{tools/eval\_heatmap\_postprocess.py}
        mantiene la salida \texttt{per\_query.csv}, pero añade un
        \texttt{aggregate\_by\_slice.csv} y tablas adicionales en
        \texttt{summary.md} para resumir el rendimiento del heatmap raw y del
        heatmap postprocesado por \texttt{query\_type} y por \texttt{tag}. Para
        que esta tercera opción del submenú no requiera rutas ni carpetas
        manuales, se añade \texttt{tools/run\_heatmap\_offline\_eval.py}, que
        normaliza la salida bajo
        \texttt{/workspace/results/eval\_runs/\{timestamp\}/heatmap\_offline/}
        y escribe también su propio \texttt{config.json}.

  \item \textbf{Validación rápida sin navegación larga}: se ejecutan únicamente
        comprobaciones ligeras, coherentes con el criterio del proyecto de no
        lanzar desde el asistente corridas largas de simulación sin supervisión
        del usuario. Han pasado \texttt{bash -n} sobre el menú, \texttt{py\_compile}
        de los scripts y entrypoints tocados, una batería \texttt{pytest} nueva
        con manifests sintéticos para el agregador 2x2
        (\texttt{tests/test\_eval\_2x2\_tools.py}) y, además, las suites ya
        existentes de Fase~F, Fase~G y resolución de instancias de habitación
        (26 tests). Con ello queda validado el armazón de evaluación; lo que
        falta a partir de ahora es usarlo sobre escenas y JSONLs reales para
        producir resultados y discutirlos en la memoria.

\end{enumerate}

% =============================================================================
\subsection*{23 de abril de 2026}
% =============================================================================

\begin{enumerate}

  \item \textbf{Limpieza del dataset HSSD}: se reduce el listado de escenas
        disponibles a las tres que se usarán en la evaluación oficial:
        \texttt{102344193\_0}, \texttt{102344280\_0} y
        \texttt{108736884\_177263634\_4}. Las restantes (variantes y dos
        carpetas con nombres corruptos generadas por capturas accidentales de
        \texttt{stderr} durante el \emph{collect}) se eliminan de
        \texttt{data/vlmaps\_dataset\_hssd/}. El menú interactivo refleja ya
        sólo \texttt{scene\_id} 0/1/2.

  \item \textbf{Reversión de la rama room-aware en el harness de evaluación}:
        tras un día con seis métodos
        (orquestador~$\times$~heatmap~$\times$~room-aware) se decide volver al
        2$\times$2 original (\texttt{Ob\_Hb}, \texttt{Ob\_Hp}, \texttt{Oe\_Hb},
        \texttt{Oe\_Hp}). Razón: el switch room-aware no es una variable
        independiente de evaluación sino una decisión arquitectónica del
        executor, y mantenerlo como eje del harness multiplicaba el coste sin
        aportar comparaciones nuevas para el nivel actual (muebles y
        habitaciones). \texttt{tools/eval\_methods.py} y
        \texttt{tools/run\_full\_eval.py} vuelven a la matriz de cuatro celdas.

  \item \textbf{Calibración del umbral de confianza de YOLOE}: se introduce un
        sweep dedicado al verificador visual sin barrer los cuatro métodos del
        2$\times$2. Justificación: el conf-thresh sólo afecta la fase de
        verificación al llegar al candidato; las decisiones de navegación,
        heatmap y selección de habitación son independientes. Por tanto basta
        con barrer un único método (\texttt{Oe\_Hp}) para aislar el efecto del
        umbral.

  \item \textbf{Nueva herramienta \texttt{tools/run\_yoloe\_thresh\_sweep.py}}
        (encargada a Codex). Corre \texttt{run\_nav\_eval.py} repetidamente
        sobre la misma escena y JSONL, una vez por threshold, propagando el
        valor vía \texttt{VLMAPS\_YOLOE\_CONF\_THRESH}. Al terminar, agrega los
        manifests y produce \texttt{compare\_thresh.csv} y
        \texttt{compare\_thresh.md} con una fila por threshold y columnas SR,
        Object~SR, Found~rate, Wrong~visits, CFR, CT2R, Early~Stop~Rate.
        Convive sin tocar el harness 2$\times$2 ni el submódulo \texttt{vlmaps}.

  \item \textbf{Resultado del barrido sobre $\{0.30, 0.40, 0.50, 0.60\}$}
        (escena \texttt{102344193\_0}, 50 queries, método \texttt{Oe\_Hp}):
        $\tau=0.30$ maximiza SR (\textbf{0.82}), Object~SR (\textbf{0.76}) y
        \emph{found rate} (\textbf{0.62}) y minimiza las visitas erróneas
        (\textbf{0.84}). A partir de $\tau=0.40$ la degradación es monótona:
        las visitas erróneas se triplican y la SR cae 16~puntos en $\tau=0.60$.
        La constancia de CFR ($0.18$ en los cuatro runs) confirma que el
        experimento aísla correctamente el verificador y no contamina la
        navegación. Los detalles cuantitativos y la decisión están integrados
        en la subsección \emph{Calibración del umbral de confianza de YOLOE}
        del bloque de testing.

  \item \textbf{Discusión honesta sobre la medición de falsos positivos}:
        durante el análisis se descarta el indicador \emph{FP-proxy} inicial
        (\texttt{found} $\land$ $\neg$\texttt{success} sobre queries
        \texttt{room\_object}). Razones: (i)~el \emph{ground truth} del
        \emph{benchmark} es a nivel de habitación, no de instancia, por lo que
        encontrar una instancia válida del mismo objeto en otra habitación se
        contabiliza erróneamente como falso positivo; (ii)~el agente detiene la
        búsqueda en la primera confirmación, ocultando posibles FPs latentes en
        otras habitaciones; (iii)~con sólo 8 queries \texttt{room\_object} el
        indicador es estadísticamente inservible. Se documenta como limitación
        del banco de evaluación que una medición directa de FPs requeriría
        \emph{queries negativas} (objetivo ausente en la escena) o
        \emph{ground-truth} de instancia.

  \item \textbf{Decisión adoptada}: se mantiene $\tau=0.30$ como umbral de
        confianza para todos los métodos del 2$\times$2 oficial. La elección
        queda justificada documentalmente (tabla, gráficos y discusión) en la
        sección de testing, evitando aparecer como un default arbitrario.

  \item \textbf{Decisión estratégica: activar la migración a ROS\,1 Noetic
        antes que el nivel~3}. Tras revisar el coste real de añadir nivel~3
        (objetos pequeños) directamente sobre Habitat y luego portarlo a ROS,
        se concluye que esa secuencia obliga a iterar dos veces sobre los
        mismos componentes (Phase~G, surrogate furniture, YOLOE bridge,
        harness). Se reordena el roadmap para migrar primero el stack a
        ROS\,1 Noetic + Gazebo y desarrollar el nivel~3 directamente sobre el
        stack final. Plazo asumido para la migración completa: \textbf{un mes}
        (hasta el 23~de~mayo de 2026).

  \item \textbf{Renuncia explícita a HSSD como dependencia obligatoria}: las
        escenas pasarán a ser \emph{house worlds} de Gazebo (HSR oficial,
        AWS RoboMaker) o iGibson (15 apartamentos basados en escaneos reales,
        utilizado por MoMa-LLM). El razonamiento por habitaciones, surrogate
        furniture y verificación con YOLOE son agnósticos a la fuente de la
        escena.

  \item \textbf{Esquema de etiquetado de habitaciones por centroides
        + Voronoi} (estilo MoMa-LLM): se sustituye el sistema de polígonos
        \texttt{semantic\_config.json} por una partición de Voronoi sobre
        centroides marcados manualmente (un \emph{click} por habitación). El
        mismo esquema vale para escenas Gazebo, iGibson y, más adelante, para
        el laboratorio físico. Coste de anotación: 5--10 minutos por escena.
        Implementación en \texttt{room\_provider} estimada en una sesión
        manteniendo la interfaz pública.

  \item \textbf{Inspiración explícita en MoMa-LLM} (Honerkamp et~al.,
        RSS SemRob 2024): adopción conceptual de tres ideas del paper, con
        re-implementación propia y \emph{sin copiar código} (compromiso de
        integridad académica). Las ideas adoptadas son: (i)~clustering de
        habitaciones por Voronoi; (ii)~vocabulario de acciones de alto nivel
        (ya parcialmente cubierto por la Phase~F existente); (iii)~métrica
        AUC-E (área bajo la curva de eficiencia) como complemento al SR a
        \emph{budget} fijo, para reportar resultados en la memoria.

  \item \textbf{Continuidad de Docker en host distinto}: la migración mantiene
        Docker como mecanismo de empaquetado. El contenedor pasa de
        \texttt{tfg-sim} (Habitat-Sim sobre Ubuntu~20.04) a un nuevo contenedor
        Noetic+Gazebo sobre Ubuntu~20.04. El usuario sigue trabajando desde su
        host Ubuntu~22.04 sin reinstalar nada en el host. Cuando se traslade el
        proyecto a un equipo nativo Ubuntu~20.04, el contenedor sigue siendo
        válido (ROS nativo en el host se considera optimización opcional, no
        requisito).

\end{enumerate}

\subsection*{24 de abril de 2026 --- Inserción real de objetos pequeños y \emph{deconfliction} del placer}

\begin{enumerate}
  \item \textbf{Reducción del dataset HSSD a tres escenas operativas}:
        tras la limpieza del 23.04, las escenas activas para evaluación
        de objetos pequeños quedan fijadas en \texttt{102344193\_0}
        (single-floor, 8 habitaciones, balcony interior), \texttt{102344280\_0}
        (15 habitaciones, ${\sim}447\,\text{m}^2$, escena principal del
        proyecto) y \texttt{108736884\_177263634\_4} (escena con
        \texttt{office} explícita, necesaria para validar el caso
        multi-room). Otras escenas evaluadas
        (\texttt{107734254}, \texttt{107734479} y variantes) se
        descartan: bancos y otomanas bloqueando el piano impedían el
        plan de movimiento del robot, y la cobertura del navmesh era
        irregular tras los \emph{collect}. Las carpetas se eliminan de
        \texttt{data/vlmaps\_dataset\_hssd/} desde dentro del contenedor
        para evitar conflictos de permisos.
  \item \textbf{Spawn fijo por escena en \texttt{collect\_hssd\_dataset.py}}:
        se añade el flag \texttt{-{}-start-pos X Z} al script de
        colección. El valor se snappea al navmesh con
        \texttt{sim.pathfinder.snap\_point(\dots)} y la \texttt{Y} se
        deja libre para que la lleve la malla. El menú interactivo
        (\texttt{docker/menu.sh}) define un diccionario
        \texttt{\_SPAWN\_POINTS} por \texttt{scene\_id} y propaga el
        flag automáticamente, evitando \emph{spawns} aleatorios encima
        de camas o dentro de armarios que invalidaban la colección.
  \item \textbf{Inventario de objetos a colocar}: se redacta una primera
        \emph{spec} JSON por escena bajo
        \texttt{tools/eval\_queries/specs/small\_objects\_*.json} con
        ocho entradas (\texttt{bottle}, \texttt{laptop}, \texttt{mug},
        \texttt{kettle}, \texttt{toaster}, \texttt{coffee\,maker},
        \texttt{teapot}, \texttt{trash\,bin}). Cada entrada lleva
        \texttt{furniture}, \texttt{room\_hint} y \texttt{count}. Se
        crea \texttt{tools/place\_small\_objects.py} que carga la
        \emph{spec}, busca instancias de la \texttt{furniture} objetivo
        dentro de la habitación indicada y escribe un nuevo
        \texttt{*.scene\_instance.json} con los \emph{placements}
        materializados. El fichero original se preserva como
        \texttt{*.before\_placements.json}.
  \item \textbf{Ampliación de \texttt{room\_priors.py} y
        \texttt{object\_priors.py} con los nuevos targets nivel-3}: se
        añaden entradas para \texttt{kettle}, \texttt{toaster},
        \texttt{coffee maker}, \texttt{mug}, \texttt{trash bin},
        \texttt{printer}, \texttt{basketball}, etc., tanto en la
        \texttt{\_MANUAL\_TABLE} como en la tabla de
        co-ocurrencia y en los \texttt{\_SEMANTIC\_ALIASES}. El cambio
        es aditivo: ningún prior existente se modifica.
  \item \textbf{Diagnóstico del bug de placement (\texttt{kettle 0/6})}:
        el primer \emph{full eval} con la spec inicial (50 queries)
        produce SR global 14\% (\emph{baseline}). Tras varias rondas
        de afinado de la pipeline de búsqueda (\texttt{post\_fix\_v1}
        $\rightarrow$ \texttt{v7}) se llega a SR 44\% pero
        \texttt{kettle} sigue dando \textbf{0/6} en todos los runs.
        Inspeccionando el \texttt{scene\_instance.json} resultante se
        descubre que \emph{los seis \emph{targets} de cocina}
        (\texttt{bottle}, \texttt{kettle}, \texttt{toaster},
        \texttt{coffee\,maker}, \texttt{mug}, \texttt{teapot})
        comparten \emph{exactamente las mismas coordenadas XYZ} sobre
        el primer \texttt{counter} del \texttt{kitchen}: el selector
        determinista del placer devolvía siempre el primer candidato y
        ninguna variabilidad espacial. Habitat colapsa los seis objetos
        en el mismo punto y YOLOE solo ve uno.
  \item \textbf{Fix: \emph{deconfliction} por (template\_id, XZ) con
        \emph{grid offset}}: el placer se modifica para mantener un
        registro \texttt{used\_host\_keys} y un contador
        \texttt{host\_reuse\_count}. Antes de aceptar un candidato, se
        reordena la lista de \emph{candidates} de manera
        \emph{fresh-first}: hosts no usados se prueban antes de reciclar
        uno ya ocupado. Cuando la reutilización es inevitable, se aplica
        un \emph{offset} determinista en una rejilla 3$\times$3
        (\texttt{\_grid\_offset}) con paso de \textbf{35\,cm} sobre
        \emph{furniture} y de \textbf{30\,cm} sobre suelo. El cambio es
        puramente aditivo y no toca la lógica de búsqueda. La
        verificación a mano sobre \texttt{102344193\_0} muestra los
        seis objetos repartidos por el counter en la rejilla esperada.
  \item \textbf{Resultado del \emph{run} \texttt{post\_placement\_full}}:
        SR global pasa de 44\% a \textbf{68\%} (34/50). El salto de
        +24~puntos viene de un único cambio no relacionado con la
        búsqueda. \texttt{kettle} pasa de 0/6 a 3/6 y \texttt{teapot}
        de 4/8 a 7/8. Las regresiones (\texttt{trash\,bin} 3/6 $\to$ 2/6)
        se atribuyen al \emph{floor offset} de 80\,cm que mueve el bin
        fuera del centroide de la habitación; queda como TODO reducirlo.
  \item \textbf{Bug de routing identificado pero no corregido}: la query
        \emph{``the laptop in the bedroom''} se enruta a la sala de estar
        cuando los \emph{room\_scores} están casi empatados
        (\texttt{bedroom}=0.37 vs \texttt{living\,room}=0.35). El
        \emph{room hint} explícito de la query no rompe el empate. Es un
        bug del selector de habitación, fuera del alcance del placer; se
        registra en \texttt{tools/HOWTO\_validate\_small\_objects.md}
        como follow-up para
        \texttt{vlmaps/policy/strategic\_policy.py}.
\end{enumerate}

\subsection*{27 de abril de 2026 --- Desbloqueo de \texttt{laptop} en el catálogo y criba del set de objetos}

\begin{enumerate}
  \item \textbf{Diagnóstico del bug de catálogo de \texttt{laptop}}:
        el placer registraba sistemáticamente
        \texttt{[skip] no HSSD template for object 'laptop'} en dos de
        las tres escenas. El CSV
        \texttt{object\_categories\_filtered.csv} contenía únicamente
        identificadores de la familia \texttt{Laptop\_X}, que no llevan
        \texttt{.object\_config.json} en HSSD. El filtro
        \texttt{\_available\_object\_handles()} los descartaba a todos
        antes de llegar a la fase de elección, dejando la categoría
        vacía.
  \item \textbf{Fix: \texttt{\_EXTRA\_CATALOG} aditivo}: se añade un
        diccionario \texttt{\_EXTRA\_CATALOG} en
        \texttt{tools/place\_small\_objects.py} con ocho hashes de
        portátiles reales (Toshiba Satellite, intel celeron 15.4,
        Computer Laptop, Sony Vaio, HP Pavilion, Toshiba L300D, MacBook
        Pro 17 y MacBook Pro 15.4) que sí traen el \texttt{.object\_config.json}
        completo. \texttt{\_load\_object\_catalog()} fusiona estos
        hashes a posteriori, respetando el filtro de disponibilidad
        en disco. El cambio es aditivo: el resto del catálogo CSV se
        carga igual.
  \item \textbf{Reducción del \emph{floor offset} a 30\,cm}: el offset
        para \emph{placements} sobre suelo (usado por \texttt{trash bin})
        se baja de 80\,cm a 30\,cm. Ya no hay regresión teórica para el
        bin actual y se protege contra futuros \emph{multi-floor placements}.
  \item \textbf{Criba del catálogo de objetos pequeños}: tras revisar
        los resultados de \texttt{post\_placement\_full} y discutir las
        propiedades semánticas de cada \emph{target}, se decide
        \emph{quitar} \texttt{trash bin} y \texttt{kettle} y
        \emph{añadir} \texttt{book}. El razonamiento:
        \begin{itemize}
          \item \texttt{trash bin} aparece en cinco habitaciones
                (\texttt{kitchen}, \texttt{bathroom}, \texttt{bedroom},
                \texttt{office}, \texttt{hallway}) sin prior dominante.
                Esto vuelve la búsqueda \emph{room-aware} indistinguible
                de una búsqueda plana M2: el sistema acaba probando
                casi todas las habitaciones. Mide ruido del detector,
                no calidad del razonamiento jerárquico.
          \item \texttt{kettle} comparte exactamente el mismo prior que
                \texttt{teapot} (\texttt{kitchen}\,$=$\,1.0), morfología
                similar (asa $+$ pico) y YOLOE los confunde
                sistemáticamente. Tener los dos en la batería no añade
                señal experimental e infla la matriz de confusión.
          \item \texttt{book} aporta el caso multi-room interpretable
                (\texttt{office}\,${>}$\,\texttt{bedroom}\,${>}$\,\texttt{living}),
                forma distintiva, abundancia de \emph{assets} con
                \texttt{.object\_config.json} en HSSD y prior ya
                cargado en \texttt{room\_priors.py}. Es justo el tipo de
                consulta donde el razonamiento jerárquico debe destacar
                frente al M2 plano.
        \end{itemize}
  \item \textbf{\emph{Specs} actualizadas}: las tres
        \texttt{tools/eval\_queries/specs/small\_objects\_*.json}
        pasan de 8 a 7 placements (50~$\to$~45 queries efectivas).
        Distribución del nuevo \texttt{book}: en
        \texttt{102344193\_0} y \texttt{102344280\_0} sobre
        \texttt{table} en \texttt{bedroom}; en
        \texttt{108736884\_177263634\_4} sobre \texttt{table} en
        \texttt{living room} para diferenciarlo del \texttt{laptop} que
        ya está en \texttt{office}.
  \item \textbf{Documentación operativa}: se redacta
        \texttt{tools/HOWTO\_validate\_small\_objects.md} con tres
        procedimientos progresivos (\emph{quick check} de 1\,min, smoke
        test de 5\,min, \emph{full eval} de ${\sim}$30\,min) y los
        criterios de fallo esperados, de modo que la validación
        post-placement se pueda repetir sin asistencia.
  \item \textbf{Aclaración metodológica para la fase posterior}: el
        régimen actual coloca cada objeto en su zona más probable
        (caso \emph{happy path} room-aware). La métrica que cierra la
        hipótesis exigirá generar \emph{N}~variantes del mismo mapa con
        colocaciones plausibles distintas (botella en encimera vs.\
        comedor vs.\ mesilla, libro en escritorio vs.\ mesilla vs.\
        sofá) y comparar cómo degradan el sistema room-aware y el M2
        plano cuando el objeto no está en su prior. Esa generación
        automática y la agregación cross-variant quedan como trabajo
        siguiente a \texttt{post\_placement\_v2}.
\end{enumerate}

\subsection*{28 de abril de 2026 --- Separación estable HSSD/ROS y arranque del Toyota HSR oficial en Gazebo}

\begin{enumerate}
  \item \textbf{Separación explícita de los dos frentes del proyecto}: se fija
        que \texttt{tfg-sim} conserva intacto el pipeline Habitat/HSSD
        (colección, evaluación, queries y métricas SR/CFR), mientras que
        \texttt{tfg-ros} concentra la migración ROS\,1 Noetic + Gazebo. Esta
        separación evita que los experimentos HSSD se rompan al iterar sobre
        Gazebo, RViz, sensores o modelos del robot.
  \item \textbf{Dos perfiles de robot en ROS}: se mantiene
        \texttt{hsr\_proxy\_gazebo\_move\_base.launch} como backend navegable
        ligero para validar \texttt{move\_base}, \texttt{rosbridge} y
        \texttt{/vlmap/navigation\_result}. En paralelo se añade
        \texttt{hsr\_gazebo\_move\_base.launch} para lanzar el Toyota HSR
        oficial público como \texttt{hsrb}, con mallas, árbol TF y sensores más
        fieles.
  \item \textbf{Modelos oficiales públicos sin vendorear assets}: se incorporan
        localmente \texttt{hsr\_meshes} y \texttt{hsr\_description} bajo
        \texttt{ros\_ws/src/external/}. La reproducción queda automatizada con
        \texttt{tools/fetch\_hsr\_official\_models.py}, pero los repositorios
        externos se excluyen de Git para no subir mallas ni documentación de
        terceros al repositorio del TFG.
  \item \textbf{Compatibilidad Noetic del Xacro oficial}: el modelo público
        requería adaptar sintaxis antigua de Xacro
        (\texttt{macro} $\rightarrow$ \texttt{xacro:macro},
        \texttt{insert\_block} $\rightarrow$ \texttt{xacro:insert\_block}) y
        corregir nombres de transmisiones generados. El síntoma visible era que
        el HSR aparecía sin cabeza o con enlaces colapsados en el origen; tras
        los parches la cabeza queda situada correctamente y el modelo se puede
        inspeccionar en Gazebo.
  \item \textbf{Topics HSR alineados con sensores reales}: la configuración
        ROS ya referencia \texttt{/hsrb/base\_scan},
        \texttt{/hsrb/odom\_ground\_truth},
        \texttt{/hsrb/head\_rgbd\_sensor/rgb/image\_rect\_color},
        \texttt{/hsrb/head\_rgbd\_sensor/depth\_registered/image\_rect\_raw},
        \texttt{/hsrb/head\_rgbd\_sensor/rgb/camera\_info},
        \texttt{/hsrb/head\_rgbd\_sensor/depth\_registered/camera\_info} y
        \texttt{/hsrb/head\_rgbd\_sensor/depth\_registered/rectified\_points}.
        El siguiente checkpoint técnico no es mover todavía el robot físico,
        sino visualizar y validar estos topics en RViz.
  \item \textbf{Próximos pasos}: primero se debe construir un panel RViz con
        \texttt{RobotModel}, \texttt{TF}, laser, RGB, depth y nube de puntos;
        después validar la cadena \texttt{map $\rightarrow$ odom
        $\rightarrow$ base\_footprint $\rightarrow$ base\_link
        $\rightarrow$ head\_rgbd\_sensor\_rgb\_frame}; después conectar
        \texttt{yoloe\_verifier} a un topic RGB configurable; después resolver
        la movilidad del HSR oficial mediante controladores TMC o mantener el
        proxy como backend de navegación; y finalmente crear un \emph{house
        world} Gazebo con habitaciones anotadas mediante centroides + Voronoi.
        Manipulación, MoveIt, brazo y grasping se posponen hasta que navegación
        y percepción sean estables.
\end{enumerate}

\subsection*{4 de mayo de 2026 --- Consolidación del backend ROS/HSR y aceleración gráfica}

\begin{enumerate}
  \item \textbf{HSR oficial como camino principal en Gazebo}: el frente ROS deja
        de apoyarse en el robot \emph{proxy} como referencia de uso diario. Se
        estabiliza el arranque del Toyota HSR oficial en mundos simples
        (\texttt{hsr\_gazebo\_walls.launch}, \texttt{hsr\_gazebo\_simple.launch}
        y variantes indoor), manteniendo \texttt{tfg-sim} intacto como backend
        Habitat/HSSD del TFG.
  \item \textbf{Sensores visibles en RViz}: queda operativo un panel RViz capaz
        de mostrar \texttt{RobotModel}, \texttt{TF}, \texttt{LaserScan},
        RGB-D, depth, nube de puntos y marcadores auxiliares. Esto convierte la
        validación del HSR en una inspección visual reproducible, sin depender
        sólo de logs o topics aislados.
  \item \textbf{Teleoperación estándar ROS}: se sustituye el teleop casero por
        \texttt{teleop\_twist\_keyboard}, remapeado a
        \texttt{/hsrb/command\_velocity}. Con ello la teleoperación pasa a usar
        una herramienta estándar del ecosistema ROS\,1 y deja de depender de un
        script ad hoc del proyecto.
  \item \textbf{Diagnóstico y fix de rendimiento gráfico en \texttt{tfg-ros}}:
        el contenedor ROS estaba renderizando Gazebo con
        \texttt{Mesa llvmpipe} (software rendering por CPU). Se habilita acceso
        NVIDIA en \texttt{docker-compose.yml} para \texttt{tfg-ros}, con
        capacidades gráficas adecuadas, y se verifica que tanto Gazebo como
        OpenGL pasan a usar la GPU NVIDIA real. El impacto práctico es que
        Gazebo GUI y RViz dejan de estar penalizados por rasterización software.
  \item \textbf{Escenarios mínimos para validación}: además de los mundos
        internos de Gazebo, se añade un mundo mínimo con paredes
        (\texttt{minimal\_walls.world}) para depurar navegación, sensores y
        render en el caso más simple posible. Esto da un banco estable para
        comprobar cámara, lidar, odometría y teleop antes de saltar a mundos más
        complejos.
  \item \textbf{Estado al cierre del día}: la arquitectura queda ya dividida en
        dos caminos claros. \texttt{tfg-sim} conserva la lógica semántica,
        HSSD, métricas y evaluación; \texttt{tfg-ros} dispone de Gazebo, RViz,
        HSR oficial, teleop estándar y aceleración GPU funcional. El siguiente
        paso ya no es infraestructura básica, sino seguir refinando la conexión
        entre ambos contenedores y el uso del backend ROS como ejecutor de goals
        semánticos.
\end{enumerate}

\subsection*{6 de mayo de 2026 --- Regiones LabelMe y exploración activa de habitaciones en Habitat/HSSD}

\begin{enumerate}
  \item \textbf{Regiones manuales como fuente prioritaria}: el flujo HSSD acepta
        ahora mapas de habitación creados con LabelMe. Si existe
        \texttt{room\_map/} para una escena, \texttt{LabelMeRoomProvider} se
        usa por delante de las regiones nativas de HSSD. Esto permite definir
        habitaciones como zonas navegables reales, no sólo como centroides o
        polígonos semánticos imperfectos.
  \item \textbf{Escenas preparadas}: las escenas \texttt{102344193\_0} y
        \texttt{102344280\_0} ya tienen regiones LabelMe convertidas y
        cargables por \texttt{interactive\_object\_nav.py}. La escena
        \texttt{108736884\_177263634\_4} se retira del menú HSSD para evitar
        confusión con la variante \texttt{108736884\_177263634\_0}.
  \item \textbf{Menos escaneos redundantes}: se añade memoria de zonas de
        escaneo por episodio. Si el robot vuelve a una posición situada a menos
        de \texttt{VLMAPS\_SCAN\_DEDUP\_RADIUS\_M} metros de una zona ya
        escaneada, se evita repetir el barrido local. El valor por defecto es
        $1{,}5$\,m.
  \item \textbf{Exploración de habitación con YOLOE activo}: al activar
        \texttt{VLMAPS\_ROOM\_EXPLORATION=1}, si falla el candidato principal
        y fallan los muebles alternativos de la habitación, el sistema genera
        puntos de paso dentro de la región LabelMe, los ordena por proximidad y
        navega por ellos con YOLOE ejecutándose durante el movimiento. Si YOLOE
        dispara con umbral bajo, el robot pausa la ruta, centra la vista, se
        acerca y confirma con umbral alto. Si la confirmación falla, registra
        la zona como falsa alarma y continúa explorando.
  \item \textbf{Parámetros de control}: la exploración queda acotada por
        \texttt{VLMAPS\_EXPLORE\_MAX\_POINTS} (8 por defecto),
        \texttt{VLMAPS\_EXPLORE\_TIMEOUT\_S} (60\,s),
        \texttt{VLMAPS\_YOLOE\_ROOM\_LOW\_THRESH} (0,40),
        \texttt{VLMAPS\_YOLOE\_CONFIRM\_THRESH} (0,65) y
        \texttt{VLMAPS\_MAX\_APPROACH\_ATTEMPTS} (3). Por defecto la
        exploración está desactivada para no alterar experimentos previos.
  \item \textbf{Telemetría nueva}: \texttt{SearchState} y
        \texttt{[eval-summary]} incluyen
        \texttt{entered\_zone\_exploration},
        \texttt{n\_exploration\_points\_planned},
        \texttt{n\_exploration\_points\_visited},
        \texttt{n\_unique\_scan\_zones},
        \texttt{n\_low\_conf\_detections},
        \texttt{n\_approach\_attempts},
        \texttt{n\_false\_positive\_zones},
        \texttt{continuous\_yoloe\_triggered} y
        \texttt{final\_success\_source}. Estas métricas permiten separar
        búsquedas resueltas por candidato directo, por mueble alternativo o por
        exploración de la habitación.
  \item \textbf{Prueba manual desde el menú}: abrir \texttt{menu}
        $\rightarrow$ \texttt{VLMaps pipeline} $\rightarrow$ \texttt{HSSD}
        $\rightarrow$ \texttt{x) LLM navigation + room exploration}. Esa opción
        activa \texttt{VLMAPS\_ROOM\_EXPLORATION=1}, pregunta los límites de
        exploración y lanza \texttt{interactive\_object\_nav.py} sin tener que
        preparar variables manualmente. La opción \texttt{l) Interactive LLM
        navigation} mantiene el comportamiento clásico sin exploración activa.
  \item \textbf{Doble capa LabelMe/Voronoi}: se separan dos usos que antes
        estaban mezclados. \texttt{room\_map.npy} conserva solo el suelo
        navegable marcado manualmente y se usa para movimiento, puntos de paso
        y límites duros de exploración. La nueva capa
        \texttt{room\_voronoi.npy} expande cada región navegable hacia las
        celdas no etiquetadas cercanas y se usa solo para asignar muebles y
        candidatos VLMap a una habitación. Así una mesa situada sobre una celda
        ocupada puede pertenecer a \texttt{living\_room}, mientras el robot
        sigue caminando únicamente por suelo navegable.
  \item \textbf{Búsqueda confinada por habitación}: en órdenes como
        \texttt{go to the living room and find the drill}, la habitación
        explícita pasa a ser una restricción dura. El robot debe entrar en la
        zona LabelMe de \texttt{living\_room}, filtrar muebles mediante
        Voronoi y explorar solo puntos de paso dentro de la máscara navegable
        de esa habitación.
\end{enumerate}

\begin{figure}[h]
\centering
\begin{tikzpicture}[
  node distance=1.0cm and 1.1cm,
  box/.style={draw, rounded corners, thick, align=center, minimum width=3.0cm, minimum height=0.85cm},
  data/.style={box, fill=blue!8},
  proc/.style={box, fill=orange!12},
  use/.style={box, fill=green!10},
  arrow/.style={-{Latex[length=3mm]}, thick}
]
  \node[data] (labelme) {LabelMe\\suelo navegable};
  \node[proc, right=of labelme] (raster) {Rasterizar\\\texttt{room\_map.npy}};
  \node[proc, right=of raster] (voronoi) {Expandir por\\Voronoi};
  \node[data, right=of voronoi] (rv) {\texttt{room\_voronoi.npy}\\pertenencia};

  \node[use, below=1.15cm of raster] (nav) {Movimiento\\y puntos de paso};
  \node[use, below=1.15cm of rv] (furn) {Muebles VLMap\\y candidatos};

  \draw[arrow] (labelme) -- (raster);
  \draw[arrow] (raster) -- (voronoi);
  \draw[arrow] (voronoi) -- (rv);
  \draw[arrow] (raster) -- (nav);
  \draw[arrow] (rv) -- (furn);
\end{tikzpicture}
\caption{Arquitectura de doble capa: LabelMe define dónde puede pisar el robot; Voronoi define a qué habitación pertenecen los muebles y candidatos.}
\end{figure}

\begin{figure}[h]
\centering
\begin{tikzpicture}[
  node distance=0.8cm,
  state/.style={draw, rounded corners, thick, align=center, minimum width=4.1cm, minimum height=0.72cm, fill=gray!8},
  good/.style={state, fill=green!10},
  warn/.style={state, fill=orange!12},
  arrow/.style={-{Latex[length=3mm]}, thick}
]
  \node[state] (parse) {Orden: habitación + objeto};
  \node[state, below=of parse] (pin) {Fijar habitación explícita};
  \node[state, below=of pin] (enter) {Entrar en zona LabelMe navegable};
  \node[warn, below=of enter] (furniture) {Muebles candidatos\\filtrados con Voronoi};
  \node[warn, below=of furniture] (scan) {Mirar mueble / punto\\barrido visual corto};
  \node[state, below=of scan] (cover) {Cobertura de habitación\\solo dentro de LabelMe};
  \node[good, below=of cover] (end) {Confirmar objeto o fallar};

  \draw[arrow] (parse) -- (pin);
  \draw[arrow] (pin) -- (enter);
  \draw[arrow] (enter) -- (furniture);
  \draw[arrow] (furniture) -- (scan);
  \draw[arrow] (scan) -- (cover);
  \draw[arrow] (cover) -- (end);
\end{tikzpicture}
\caption{Flujo previsto para búsquedas confinadas: la habitación explícita limita tanto la exploración como los muebles inspeccionables.}
\end{figure}

\subsection*{7 de mayo de 2026 --- Corrección de LabelMe/Voronoi y validación de regiones}

\begin{enumerate}
  \item \textbf{Sincronización automática antes de navegar}: las opciones
        \texttt{l) Interactive LLM navigation} y \texttt{x) LLM navigation +
        room exploration} reconvierten automáticamente el JSON de LabelMe si
        existe en \texttt{/workspace/annotations/room\_labels}. Esto evita que
        el sistema cargue un \texttt{room\_map} antiguo cuando se han añadido
        o renombrado regiones desde LabelMe.
  \item \textbf{Validación de etiquetas}: al arrancar con
        \texttt{VLMAPS\_ROOM\_EXPLORATION=1}, el sistema compara las etiquetas
        del JSON de LabelMe con las habitaciones cargadas por
        \texttt{LabelMeRoomProvider}. Si falta alguna región, aborta con un
        mensaje explícito en vez de explorar con un mapa incompleto.
  \item \textbf{Voronoi limitado por paredes}: se sustituye la expansión
        euclídea global por una expansión conectada desde las regiones LabelMe
        dentro de un dominio interior visible. El dominio se obtiene de
        \texttt{topdown\_labeled.png}, \texttt{topdown\_rgb.png} y
        \texttt{obstacle\_map.png}, conserva solo componentes conectados que
        tocan alguna región LabelMe y no permite asignar celdas fuera de ese
        dominio. La capa resultante sigue siendo solo de pertenencia semántica;
        el robot continúa caminando únicamente por \texttt{room\_map.npy}.
  \item \textbf{Artefactos de depuración}: cada conversión guarda
        \texttt{room\_voronoi\_domain.png} y
        \texttt{room\_voronoi\_viz.png}; además, el arranque de \texttt{x}
        guarda \texttt{debug/room\_debug\_overlay.png}. Estas imágenes permiten
        comprobar visualmente si la expansión se queda dentro de la vivienda y
        si los muebles caen bajo la habitación correcta.
  \item \textbf{Validación realizada}: se regeneraron las tres escenas HSSD
        etiquetadas. Resultado: \texttt{102344193\_0} carga 7/7 etiquetas,
        \texttt{102344280\_0} carga 12/12 y
        \texttt{108736884\_177263634\_0} carga 15/15. En las tres escenas,
        \texttt{outside\_domain=0}: no hay celdas Voronoi asignadas fuera del
        dominio interior generado.
\end{enumerate}

\begin{figure}[h]
\centering
\begin{tikzpicture}[
  node distance=0.85cm and 1.15cm,
  box/.style={draw, rounded corners, thick, align=center, minimum width=3.2cm, minimum height=0.78cm},
  source/.style={box, fill=blue!8},
  proc/.style={box, fill=orange!12},
  safe/.style={box, fill=green!10},
  stop/.style={box, fill=red!8},
  arrow/.style={-{Latex[length=3mm]}, thick}
]
  \node[source] (json) {JSON LabelMe\\regiones navegables};
  \node[proc, right=of json] (domain) {Dominio interior\\visible y conectado};
  \node[proc, right=of domain] (bfs) {Expansión por celdas\\máx. 50};
  \node[safe, below=of bfs] (vor) {\texttt{room\_voronoi.npy}\\muebles y candidatos};
  \node[stop, below=of domain] (walls) {Paredes / exterior\\no propagables};
  \node[safe, below=of json] (nav) {\texttt{room\_map.npy}\\movimiento};

  \draw[arrow] (json) -- (domain);
  \draw[arrow] (domain) -- (bfs);
  \draw[arrow] (bfs) -- (vor);
  \draw[arrow] (json) -- (nav);
  \draw[arrow] (walls) -- node[right, font=\footnotesize] {bloquea} (domain);
\end{tikzpicture}
\caption{Voronoi corregido: las regiones LabelMe se expanden solo por el dominio interior conectado, de modo que la pertenencia de muebles no atraviesa paredes ni exterior.}
\end{figure}

\end{document}
